% ==========================================================================
% Vorlesungsmitschrift: Grundlagen der Mathematik I
% Mathematisches Institut – Wintersemester 2026/2027
% ==========================================================================
%!TEX program = lualatex
\documentclass[11pt, a4paper]{article}

% Zentrale Konfiguration und Pakete laden
% ==========================================================================
% Präambel: Pakete, Schriftarten, Layout und Theoreme
% Angelehnt an das Design der Vorlesungsmitschrift Lineare Algebra 1
% ==========================================================================

\usepackage[a4paper, margin=2.2cm]{geometry}
\usepackage{fontspec}

% Sprach- und Schrifteinstellungen (geeignet für LuaLaTeX & XeLaTeX)
\usepackage[english, ngerman]{babel}
\babelfont{rm}{Noto Sans}

% Mathematische Pakete
\usepackage{amsmath, amssymb, amsthm}
\usepackage{booktabs}
\usepackage{enumitem}
\usepackage{xcolor}
\usepackage{hyperref}
\usepackage{tikz}
\usetikzlibrary{arrows.meta, positioning}
\allowdisplaybreaks[1]

% Farbpalette
\definecolor{primary}{RGB}{20, 50, 110}
\definecolor{secondary}{RGB}{90, 105, 125}
\definecolor{boxbg}{RGB}{246, 248, 252}
\definecolor{bordercol}{RGB}{210, 220, 235}

% Theorem-Definitionen
\theoremstyle{definition}
\newtheorem{definition}{Definition}[section]
\newtheorem{beispiel}[definition]{Beispiel}
\newtheorem{bemerkung}[definition]{Bemerkung}

\theoremstyle{plain}
\newtheorem{satz}[definition]{Satz}
\newtheorem{lemma}[definition]{Lemma}

% Listen-Formatierung
\setlist[itemize]{label=--, leftmargin=1.8em, itemsep=2pt}
\setlist[enumerate]{leftmargin=1.8em, itemsep=2pt}

% Hyperlink-Stile
\hypersetup{
    colorlinks=true,
    linkcolor=primary,
    urlcolor=primary,
    pdftitle={Grundlagen der Mathematik I - Vorlesungsmitschrift},
    pdfauthor={Soroush}
}

% Standard-Mengensymbole als handliche Makros
\newcommand{\N}{\mathbb{N}}
\newcommand{\Z}{\mathbb{Z}}
\newcommand{\Q}{\mathbb{Q}}
\newcommand{\R}{\mathbb{R}}
\newcommand{\C}{\mathbb{C}}
\newcommand{\K}{\mathbb{K}}
\newcommand{\F}{\mathbb{F}}
\newcommand{\Pot}{\mathcal{P}}
\newcommand{\xor}{\mathbin{\dot{\lor}}}

% Vorlesungs-Marker: Datum und zugehörige handschriftliche Quelle
\newcommand{\vorlesungsdatum}[2]{%
    \vspace{0.6em}%
    \noindent\textcolor{secondary}{\small\textbf{Vorlesung vom #1} \hfill \textit{Handschrift: \texttt{\detokenize{#2}}}}\par%
    \vspace{0.4em}%
}

% Seiten beziehen sich stets auf die PDF-Seiten der kumulativen Mitschrift.
\newcommand{\notizquelle}[1]{%
    \par\smallskip\noindent{\small\textcolor{secondary}{%
    Quelle: \texttt{handschrift/09112026.PDF}, PDF-S.~#1.}}\par\smallskip
}
\newcommand{\notizquelleNeu}[1]{%
    \par\smallskip\noindent{\small\textcolor{secondary}{%
    Quelle: \texttt{handschrift/Vorlesung/vorlesungsmitschrift\_gdm1PDF\_261114\_161111.PDF}, PDF-S.~#1.}}\par\smallskip
}


\begin{document}

% --------------------------------------------------------------------------
% Titelkopf
% --------------------------------------------------------------------------
\begin{center}
    {\Large\textbf{\textcolor{primary}{Grundlagen der Mathematik I}}}\\[0.4em]
    {\large fb12 -- Wintersemester 2026/2027}\\[0.2em]
    {\small Vorlesungsmitschrift -- Stand der Quellen: 17.11.2026}\\[0.4em]
    \rule{\textwidth}{0.8pt}
\end{center}

\vspace{-0.5em}

% --------------------------------------------------------------------------
% Inhaltsverzeichnis (optional ein-/ausblendbar)
% --------------------------------------------------------------------------
\noindent Grundlage sind die eigenen Mitschriften
(\texttt{09112026.PDF}, 54~PDF-Seiten, der Nachtrag
\texttt{vorlesungsmitschrift\_gdm1PDF\_261114\_161111.PDF}, 27~PDF-Seiten,
sowie \texttt{Uebungsblatt 3\_261117\_094748.PDF}, 11~PDF-Seiten)
und das vorliegende Begleitskript von Prof.~Dr.~Oliver Goertsches.
Die Quellenhinweise nennen PDF-Seiten, keine nachträglich vermuteten
Vorlesungstermine. Beispiele und Beweise sind sprachlich ausgearbeitet;
fachliche Präzisierungen werden kenntlich gemacht.
\tableofcontents
\clearpage

% --------------------------------------------------------------------------
% Kapitel einbinden
% --------------------------------------------------------------------------
% ==========================================================================
% Organisatorisches & Übungsbetrieb
% ==========================================================================

\section*{Organisatorisches \& Übungsbetrieb}
\notizquelle{1}

\begin{itemize}
    \item \textbf{Lehrveranstaltung:} Grundlagen der Mathematik I (GdM 1), Wintersemester 2026/2027.
    \item \textbf{Dozent:} Prof.~Dr.~Oliver Goertsches (Philipps-Universität Marburg, Fachbereich 12).
    \item \textbf{Übungsblätter:} Handschriftliche Abgabe jeweils \textbf{dienstags}.
    \item \textbf{Abgabeform:} Handschriftlich (Gruppenarbeit zu zweit möglich und empfohlen).
    \item \textbf{Präsenztests:} In den Übungsgruppen im Rhythmus von ca.~alle 4 Wochen.
    \item \textbf{Klausurzulassung:} 
    \begin{itemize}
        \item Mindestens \textbf{50\,\%} der Gesamtpunkte aus Übungsaufgaben und Präsenztests.
        \item \textbf{Wichtig:} Rechtzeitige Anmeldung zur Studienleistung!
    \end{itemize}
    \item \textbf{Klausurvorbereitung:} Altklausuren stehen online zur Verfügung.
\end{itemize}

\vspace{0.5em}
\hrule
\vspace{0.8em}

% ==========================================================================
% Kapitel 1: Logik
% ==========================================================================

\section{Logik}

\notizquelle{1--17}

In diesem Kapitel besprechen wir kurz und prägnant die aussagen- und prädikatenlogischen Grundlagen, die nötig sind, um Mathematik formal und widerspruchsfrei zu betreiben. Die Logik dient uns hierbei vor allem als universelle Sprache zur Formulierung von Definitionen, Sätzen und Beweisen.

% --------------------------------------------------------------------------
% Aussagen
% --------------------------------------------------------------------------
\subsection{Aussagen}

Ausgangspunkt der mathematischen Logik ist der Begriff der Aussage:

\begin{definition}[Aussage]
Eine \textbf{Aussage} ist ein sprachlicher, grammatikalisch korrekter Ausdruck (Satz), welcher eindeutig \textbf{wahr} ($1$ bzw.\ $w$) oder eindeutig \textbf{falsch} ($0$ bzw.\ $f$) ist.
\end{definition}

Die klassische Logik beruht auf dem fundamentalen \textbf{Prinzip der Zweiwertigkeit}, welches sich in zwei Kernprinzipien gliedert:
\begin{enumerate}
    \item \textbf{Prinzip vom ausgeschlossenen Widerspruch:} Eine Aussage kann nicht zugleich wahr und falsch sein.
    \item \textbf{Satz vom ausgeschlossenen Dritten (\textit{Tertium non datur}):} Eine Aussage hat stets genau einen der beiden Wahrheitswerte \glqq wahr\grqq\ oder \glqq falsch\grqq; Zwischenwerte wie \glqq vielleicht\grqq\ existieren nicht.
\end{enumerate}

\begin{bemerkung}[Abgrenzung zu Nicht-Aussagen]
Nicht jeder sprachlich wohlgeformte Satz stellt eine mathematische Aussage dar:
\begin{itemize}
    \item Fragen (\glqq Wie spät ist es?\grqq)
    \item Meinungsäußerungen und Werturteile (\glqq Dieses Bild ist schön.\grqq)
    \item Aufforderungen und Befehle (\glqq Schließe bitte die Tür!\grqq)
    \item Terme ohne Relationszeichen: $9 + 3$ ist keine Aussage (sondern ein Term); die Gleichung $9 + 3 = 12$ hingegen ist eine wahre Aussage.
    \item Paradoxien wie das Lügner-Paradoxon (\glqq Dieser Satz ist falsch.\grqq).
\end{itemize}
\end{bemerkung}

\begin{beispiel}[Wahrheitswert und Prüfbarkeit]
Ob ein Ausdruck eine Aussage ist, hängt nicht davon ab, ob der Wahrheitswert unmittelbar auf einen Blick ersichtlich ist:
\begin{itemize}
    \item \glqq $31\,274\,567$ ist eine Primzahl.\grqq\ ist eine Aussage (sie ist objektiv wahr oder falsch, auch wenn die Prüfung Rechenaufwand erfordert).
    \item \glqq Die Erde ist rund.\grqq\ $\rightarrow$ wahre Aussage.
    \item Der Satz \glqq Dieser Satz besteht aus sechs Wörtern\grqq\ enthält isoliert tatsächlich sechs Wörter. Beim Einbetten in einen längeren Satz kann sich der Bezug von \glqq dieser Satz\grqq\ jedoch ändern. Selbstreferenz ist nicht automatisch widersprüchlich, verlangt aber einen eindeutig festgelegten Kontext.
\end{itemize}
\end{beispiel}

\subsection*{Das Extensionalitätsprinzip}
\textbf{Extensionalitätsprinzip (Wahrheitsfunktionalität):} Der Wahrheitswert einer zusammengesetzten Aussage hängt ausschließlich von den Wahrheitswerten ihrer Teilkomponenten ab -- nicht von inhaltlichen, zeitlichen oder kausalen Zusammenhängen.

\begin{bemerkung}[Unterschied zur Umgangssprache]
In der Umgangssprache transportiert die Satzreihenfolge bei \glqq und\grqq\ häufig Kausalität oder zeitliche Abfolge:
\begin{enumerate}
    \item \glqq Sie verließ das Land und wurde reich.\grqq
    \item \glqq Sie wurde reich und verließ das Land.\grqq
\end{enumerate}
In der formalen Logik ist die Verknüpfung $A \land B$ rein wahrheitswertbasiert und kommutativ ($A \land B \equiv B \land A$); Zeit- und Kausalbezüge spielen keine Rolle.
\end{bemerkung}

% --------------------------------------------------------------------------
% Verknuepfungen und Wahrheitstafeln
% --------------------------------------------------------------------------
\subsection{Verknüpfungen und Wahrheitstafeln}

Seien $A$ und $B$ Aussagen. Neue, zusammengesetzte Aussagen werden über \textbf{Wahrheitstafeln} eindeutig festgelegt:

\begin{definition}[Logische Verknüpfungen]
Seien $A$ und $B$ Aussagen. Wir definieren:
\begin{enumerate}
    \item \textbf{Konjunktion ($A \land B$, \glqq $A$ und $B$\grqq):} Nur dann wahr ($1$), wenn sowohl $A$ als auch $B$ wahr sind.
    \item \textbf{Disjunktion ($A \lor B$, \glqq $A$ oder $B$\grqq):} Inklusives Oder; wahr, sobald mindestens eine der Aussagen $A$ oder $B$ wahr ist.
    \item \textbf{Antivalenz ($A \xor B$, \glqq Entweder $A$ oder $B$\grqq, XOR):} Genau dann wahr, wenn $A$ und $B$ unterschiedliche Wahrheitswerte haben.
    \item \textbf{Negation ($\neg A$ bzw.\ $\bar{A}$, \glqq Nicht $A$\grqq):} Kehrt den Wahrheitswert der Aussage um.
    \item \textbf{Implikation ($A \Rightarrow B$, \glqq Wenn $A$, dann $B$\grqq):} Nur dann falsch, wenn $A$ wahr ($1$) und $B$ falsch ($0$) ist. Ist $A$ falsch, ist die Implikation stets wahr (eine \emph{vakuos wahre Implikation}). $A$ heißt hinreichend für $B$; $B$ heißt notwendig für $A$.
    \item \textbf{Äquivalenz ($A \Leftrightarrow B$, \glqq $A$ genau dann, wenn $B$\grqq):} Wahr, wenn $A$ und $B$ denselben Wahrheitswert besitzen.
\end{enumerate}
\end{definition}

\begin{center}
\renewcommand{\arraystretch}{1.15}
\begin{tabular}{ccccccc}
    \toprule
    $A$ & $B$ & $A \land B$ & $A \lor B$ & $A \xor B$ & $A \Rightarrow B$ & $A \Leftrightarrow B$ \\
    \midrule
    0 & 0 & 0 & 0 & 0 & 1 & 1 \\
    0 & 1 & 0 & 1 & 1 & 1 & 0 \\
    1 & 0 & 0 & 1 & 1 & 0 & 0 \\
    1 & 1 & 1 & 1 & 0 & 1 & 1 \\
    \bottomrule
\end{tabular}
\qquad
\begin{tabular}{cc}
    \toprule
    $A$ & $\neg A$ \\
    \midrule
    0 & 1 \\
    1 & 0 \\
    \bottomrule
\end{tabular}
\end{center}

\begin{bemerkung}[Falsche Voraussetzung und Widerspruch]
Die Belegung von $A \Rightarrow B$ bei falscher Prämisse ($A = 0$) lässt sich an Beispielen motivieren:
\begin{itemize}
    \item \glqq Wenn ich heute im Lotto gewinne, mache ich eine Weltreise.\grqq\newline
    Gewinne ich nicht, habe ich mein Versprechen nicht gebrochen -- die Aussage bleibt wahr.
    \item Mathematisch: Die Aussagen \glqq Wenn $2+2=42$, dann ist $2+2=0$\grqq\ und \glqq Wenn $2+2=42$, dann ist die Erde eine Scheibe\grqq\ sind formal beide \textbf{wahr}, da die Voraussetzung falsch ist.
\end{itemize}
Eine falsche Voraussetzung erlaubt für sich genommen nicht, $B$ zu behaupten.
Das Prinzip \emph{ex falso quodlibet} besagt vielmehr, dass aus einem
Widerspruch jede Aussage folgt: $(A\land\neg A)\Rightarrow B$ ist eine Tautologie.
\end{bemerkung}

\subsection*{Operatorrangfolge}
In der Konvention der Vorlesung bindet die Negation stärker als
Konjunktion und Disjunktion; diese binden stärker als Implikation und
Äquivalenz. Insbesondere bedeutet $\neg A\lor B$ stets $(\neg A)\lor B$.
Bei gemischten Verknüpfungen derselben Stufe setzen wir Klammern,
z.\,B. $(A\land B)\lor C$. So wird keine zusätzliche Rangfolge vorausgesetzt.

\begin{bemerkung}[Aussage und metasprachliche Gleichwertigkeit]
$A\Leftrightarrow B$ ist eine zusammengesetzte Aussage, die bei einer
bestimmten Belegung wahr oder falsch sein kann. Dagegen besagt
$A\equiv B$, dass beide Formeln bei \emph{jeder} Belegung denselben
Wahrheitswert haben. Die Mitschrift unterscheidet dies ausdrücklich
(PDF-S.~5--7). Bei den folgenden allgemeinen Schlussregeln bezeichnet
eine hingeschriebene Äquivalenz eine Tautologie.
\end{bemerkung}

\begin{definition}[Logische Äquivalenz $\equiv$]
Zwei zusammengesetzte Aussagen $A$ und $B$ heißen \textbf{logisch äquivalent} (in Zeichen $A \equiv B$), wenn sie bei jeder Belegung der Teilkomponenten denselben Wahrheitswert liefern.
\end{definition}

\begin{satz}[Auflösung der Implikation]
Für beliebige Aussagen $A$ und $B$ gilt:
\[
(A \Rightarrow B) \;\equiv\; (\neg A \lor B).
\]
\end{satz}

\begin{proof}
Die Wahrheitstafel zeigt zeilenweise identische Ergebnisse:
\begin{center}
\renewcommand{\arraystretch}{1.15}
\begin{tabular}{ccccc}
    \toprule
    $A$ & $B$ & $\neg A$ & $A \Rightarrow B$ & $\neg A \lor B$ \\
    \midrule
    0 & 0 & 1 & 1 & 1 \\
    0 & 1 & 1 & 1 & 1 \\
    1 & 0 & 0 & 0 & 0 \\
    1 & 1 & 0 & 1 & 1 \\
    \bottomrule
\end{tabular}
\end{center}
Da die Spalten für $A \Rightarrow B$ und $\neg A \lor B$ übereinstimmen, ist die logische Äquivalenz bewiesen.
\end{proof}

% --------------------------------------------------------------------------
% Tautologien und Schlussregeln
% --------------------------------------------------------------------------
\subsection{Tautologien}

\begin{definition}[Tautologie und Widerspruch]
Eine zusammengesetzte Aussage heißt \textbf{Tautologie}, wenn sie unabhängig von den Wahrheitswerten ihrer Teilaussagen stets den Wert \glqq wahr\grqq\ ($1$) annimmt.\newline
Sie heißt \textbf{Widerspruch} (Symbol $\bot$), wenn sie stets den Wert \glqq falsch\grqq\ ($0$) annimmt.
\end{definition}

Es gilt also: $A \equiv B$ genau dann, wenn $A \Leftrightarrow B$ eine Tautologie ist.

\begin{beispiel}[Elementare Tautologien]
$A\Leftrightarrow A$, $A\lor\neg A$, $\neg(A\land\neg A)$,
$A\Rightarrow(A\lor B)$ und $(A\land B)\Rightarrow A$ sind
Tautologien. Dagegen ist $A\land\neg A$ ein Widerspruch.
\end{beispiel}

\begin{satz}[De Morgansche Regeln]
Für beliebige Aussagen $A$ und $B$ gelten:
\begin{align*}
    \neg(A \land B) &\Longleftrightarrow (\neg A \lor \neg B), \\
    \neg(A \lor B) &\Longleftrightarrow (\neg A \land \neg B).
\end{align*}
\end{satz}

\begin{proof}
Beweis der ersten Regel exemplarisch mittels Wahrheitstafel:
\begin{center}
\renewcommand{\arraystretch}{1.15}
\begin{tabular}{ccccccc}
    \toprule
    $A$ & $B$ & $\neg A$ & $\neg B$ & $A \land B$ & $\neg(A \land B)$ & $\neg A \lor \neg B$ \\
    \midrule
    0 & 0 & 1 & 1 & 0 & 1 & 1 \\
    0 & 1 & 1 & 0 & 0 & 1 & 1 \\
    1 & 0 & 0 & 1 & 0 & 1 & 1 \\
    1 & 1 & 0 & 0 & 1 & 0 & 0 \\
    \bottomrule
\end{tabular}
\end{center}
Die beiden letzten Spalten stimmen überein. Ihre Äquivalenz hat daher in jeder Zeile den Wert $1$.
\end{proof}

\begin{satz}[Klassische Schlussregeln]
Es gelten die folgenden fundamentalen Tautologien:
\begin{enumerate}
    \item \textbf{Doppelte Negation:} $A \Longleftrightarrow \neg(\neg A)$
    \item \textbf{Modus Ponens:} $(A \land (A \Rightarrow B)) \Longrightarrow B$
    \item \textbf{Gesetz des Syllogismus (Transitivität):} $((A \Rightarrow B) \land (B \Rightarrow C)) \Longrightarrow (A \Rightarrow C)$
    \item \textbf{Kontraposition (Indirekter Beweis 1):} $(A \Rightarrow B) \Longleftrightarrow (\neg B \Rightarrow \neg A)$
    \item \textbf{Widerspruchsbeweis (Indirekter Beweis 2):} $A \Longleftrightarrow (\neg A \Rightarrow \bot)$
    \item \textbf{Äquivalenzbeweis:} $(A \Leftrightarrow B) \Longleftrightarrow ((A \Rightarrow B) \land (B \Rightarrow A))$
    \item \textbf{Distributivgesetze:}
    \begin{align*}
        A \land (B \lor C) &\Longleftrightarrow (A \land B) \lor (A \land C), \\
        A \lor (B \land C) &\Longleftrightarrow (A \lor B) \land (A \lor C).
    \end{align*}
\end{enumerate}
\end{satz}

\begin{proof}[Beweis der Kontraposition ohne Wahrheitstafel]
Unter wiederholter Anwendung bereits bewiesener Äquivalenzen:
\begin{align*}
    (A\Rightarrow B)&\equiv(\neg A\lor B)\\
    &\equiv(\neg\neg B\lor\neg A)\\
    &\equiv(\neg B\Rightarrow\neg A).
\end{align*}
Dabei benutzen wir die Auflösung der Implikation, doppelte Negation und
Kommutativität der Disjunktion.
\end{proof}

\subsection*{Klassische Anwendung: Irrationalität von \texorpdfstring{$\sqrt{2}$}{sqrt(2)}}

\begin{satz}[Irrationalität von $\sqrt{2}$]
Die Zahl $\sqrt{2}$ ist irrational ($\sqrt{2} \notin \Q$).
\end{satz}

\begin{proof}
Wir führen einen indirekten Beweis (Widerspruchsbeweis): Angenommen, $\sqrt{2}$ sei rational. Dann existieren teilerfremde natürliche Zahlen $p, q \in \N_{>0}$ mit:
\[
\sqrt{2} = \frac{p}{q}.
\]
Quadrieren und Multiplikation mit $q^2$ ergibt:
\[
2q^2 = p^2.
\]
Hieraus folgt, dass $p^2$ gerade ist. Da das Quadrat einer ungeraden Zahl stets ungerade ist, muss auch $p$ gerade sein. Es existiert also ein $r \in \N$ mit $p = 2r$. Einsetzen liefert:
\[
2q^2 = (2r)^2 = 4r^2 \implies q^2 = 2r^2.
\]
Dasselbe Argument zeigt, dass auch $q$ gerade sein muss. 
Damit sind sowohl $p$ als auch $q$ durch 2 teilbar, was der Teilerfremdheit von $p$ und $q$ widerspricht. Folglich ist die Annahme falsch und $\sqrt{2}$ ist irrational.
\end{proof}

% --------------------------------------------------------------------------
% Aussageformen und Quantoren
% --------------------------------------------------------------------------
\subsection{Aussageformen und Quantoren}

\begin{definition}[Aussageform]
Eine \textbf{einstellige Aussageform} mit Einsetzungsklasse $\Omega$ ist ein sprachlicher, grammatikalisch korrekter Ausdruck $H(\cdot)$ mit einer Leerstelle, sodass für jedes $x \in \Omega$ die Einsetzung $H(x)$ eine Aussage ist.
\end{definition}

\begin{beispiel}
\begin{itemize}
    \item Sei $\Omega = \R$ und $H(x) :\Longleftrightarrow x \in \Q$. Dann ist $H(42)$ wahr, $H(\sqrt{2})$ jedoch falsch.
    \item Sei $\Omega = \R$ und $H(x) :\Longleftrightarrow x^2 \ge 0$. Dies ist für alle $x \in \R$ wahr.
    \item Sei $\Omega = \N$ und $H(n) :\Longleftrightarrow 2n+4 \text{ ist Summe zweier Primzahlen}$ (Goldbachsche Vermutung).
\end{itemize}
\end{beispiel}

\begin{definition}[All- und Existenzquantor]
Ist $H(\cdot)$ eine Aussageform über $\Omega$, so bilden wir:
\begin{itemize}
    \item $\forall x \in \Omega : H(x)$ \quad (\glqq Für alle $x \in \Omega$ gilt $H(x)$\grqq, Allquantor)
    \item $\exists x \in \Omega : H(x)$ \quad (\glqq Es gibt ein $x \in \Omega$, sodass $H(x)$ gilt\grqq, Existenzquantor)
    \item $\exists! x \in \Omega : H(x)$ \quad (\glqq Es gibt genau ein $x \in \Omega$, sodass $H(x)$ gilt\grqq)
\end{itemize}
\end{definition}

\begin{satz}[Negation von Quantoren]
Es gelten die Äquivalenzen:
\begin{align*}
    \neg(\forall x \in \Omega : H(x)) &\Longleftrightarrow \exists x \in \Omega : \neg H(x), \\
    \neg(\exists x \in \Omega : H(x)) &\Longleftrightarrow \forall x \in \Omega : \neg H(x).
\end{align*}
\end{satz}

\begin{bemerkung}[Quantoren auf Teilklassen]
Für eine Unterklasse $\Omega' \subset \Omega$ gilt:
\begin{align*}
    (\exists x \in \Omega' : H(x)) &\Longleftrightarrow \exists x \in \Omega : (x \in \Omega' \land H(x)), \\
    (\forall x \in \Omega' : H(x)) &\Longleftrightarrow \forall x \in \Omega : (x \in \Omega' \Longrightarrow H(x)).
\end{align*}
\emph{Wichtig:} Bei der Allaussage steht eine Implikation. Daher ist ein Satz wie \glqq Alle runden Dreiecke sind blau\grqq\ im mathematischen Sinne formal wahr, da die Prämisse (rundes Dreieck) niemals erfüllt ist.
\end{bemerkung}

\begin{definition}[Mehrstellige Aussageform]
Eine $n$-stellige Aussageform $H(x_1,\ldots,x_n)$ mit Einsetzungsklassen
$\Omega_1,\ldots,\Omega_n$ wird bei jeder Wahl $x_i\in\Omega_i$ zu einer
Aussage. Beispielsweise ist $H(x,y):\Longleftrightarrow x<y$ auf
$\R\times\R$ zweistellig. Ebenso ist
$H(a,b,c,d):\Longleftrightarrow a+b=c+d$ auf $\N^4$ vierstellig.
Quantoren schreiben wir vor die zugehörige Aussageform und beachten ihren
Geltungsbereich.
\end{definition}

\subsection*{Reihenfolge bei mehreren Quantoren}
Treffen unterschiedliche Quantoren aufeinander, ist die Reihenfolge entscheidend:
\begin{beispiel}
Sei $\Omega$ die Menge aller verheirateten Personen und
\[
H(x, y) :\Longleftrightarrow x \text{ ist mit } y \text{ verheiratet}.
\]
\begin{itemize}
    \item $\forall x \in \Omega \; \exists y \in \Omega : H(x, y)$:\enspace \glqq Jeder hat einen Ehepartner\grqq\ $\rightarrow$ \textbf{wahr}.
    \item $\exists y \in \Omega \; \forall x \in \Omega : H(x, y)$:\enspace \glqq Es gibt jemanden, der mit allen verheiratet ist\grqq\ $\rightarrow$ \textbf{falsch}.
\end{itemize}
\end{beispiel}
Im Ehepartnerbeispiel wird vorausgesetzt, dass niemand mit sich selbst
verheiratet ist und dass die betrachtete Menge nicht leer ist.
Ein rein mathematisches Beispiel vermeidet diese Modellannahmen:
\[
\forall x\in\R\;\exists y\in\R:\ x<y
\qquad\text{ist wahr (wähle $y=x+1$),}
\]
während $\exists y\in\R\;\forall x\in\R:\ x<y$ falsch ist
(setze für ein vorgegebenes $y$ etwa $x=y$).

\subsubsection*{Übungsblatt 3, Aufgabe 1: Quantoren, Negation und Alltagssprache}
\notizquelleDatei{handschrift/Uebungsblatt 3_261117_094748.PDF}{1--5}

\noindent\textbf{Aufgabenstellung.}
Formulieren Sie die folgenden Sätze mithilfe von Quantoren und geeignet definierten Aussagen und Aussageformen in formale Aussagen um. Benutzen Sie die Äquivalenzen
\[
    \neg\bigl(\forall x\in\Omega : H(x)\bigr) \iff \exists x\in\Omega : \neg H(x),
    \qquad
    \neg\bigl(\exists x\in\Omega : H(x)\bigr) \iff \forall x\in\Omega : \neg H(x),
\]
um die gebildeten formalen Aussagen umzuformen, und übersetzen Sie das Ergebnis wieder in alltägliche Aussagen.

\begin{enumerate}
    \item \emph{„Es gibt Blätter, die nicht auf Bäumen wachsen.“}
    \begin{itemize}
        \item \textbf{Definitionen:} Sei $B$ die Menge aller Blätter und $H(b)$ die Aussageform „$b$ wächst auf einem Baum“.
        \item \textbf{Formalisierung und Umformung:}
        \[
            \exists b\in B : \neg H(b)
            \iff \neg\bigl(\forall b\in B : H(b)\bigr).
        \]
        \item \textbf{Alltägliche Aussage:} „Es ist nicht wahr, dass alle Blätter auf Bäumen wachsen.“
    \end{itemize}

    \item \emph{„Wenn ich Eintopf koche, benutze ich alle Zutaten, die ich in der Küche finde und die noch haltbar sind.“}
    \begin{itemize}
        \item \textbf{Definitionen:} Seien $A$ die Menge aller Zutaten in meiner Küche und $B$ die Menge aller Zutaten, die noch haltbar sind. Setze $\Omega := A\cap B$, und sei $H(x)$ die Aussageform „$x$ ist in meinem selbstgekochten Eintopf“.
        \item \textbf{Formalisierung und Umformung:}
        \[
            \forall x\in\Omega : H(x)
            \iff \neg\bigl(\exists x\in\Omega : \neg H(x)\bigr)
            \iff \nexists x\in\Omega : \neg H(x).
        \]
        \item \textbf{Alltägliche Aussage:} „Es gibt keine Zutat aus meiner Küche, die noch haltbar ist und nicht in meinem Eintopf ist.“
    \end{itemize}

    \item \emph{„Für jeden handgestrickten Pullover gibt es einen Menschen, der den Pullover gestrickt hat.“}
    \begin{itemize}
        \item \textbf{Definitionen:} Sei $\Omega$ die Menge aller handgestrickten Pullover, $\Omega'$ die Menge aller Menschen und $H(x,y)$ die Aussageform „$x$ wurde von $y$ gestrickt“.
        \item \textbf{Formalisierung und Umformung:}
        \[
            \forall x\in\Omega\;\exists y\in\Omega' : H(x,y)
            \iff \neg\bigl(\exists x\in\Omega\;\forall y\in\Omega' : \neg H(x,y)\bigr).
        \]
        \item \textbf{Alltägliche Aussage:} „Es ist nicht wahr, dass es einen handgestrickten Pullover gibt, der von keinem Menschen gestrickt worden ist.“
    \end{itemize}
\end{enumerate}

\vspace{0.8em}
\hrule
\vspace{1em}

% ==========================================================================
% Kapitel 2: Mengenlehre
% ==========================================================================

\section{Mengenlehre}

\notizquelle{18--26}

Die Mengenlehre bildet die fundamentale Sprache der modernen Mathematik. Nahezu alle mathematischen Objekte -- wie Zahlen, Vektoren, Funktionen und algebraische Strukturen -- werden als Mengen oder Abbildungen zwischen Mengen modelliert.

% --------------------------------------------------------------------------
% Der Cantorsche Mengenbegriff
% --------------------------------------------------------------------------
\subsection{Der Cantorsche Mengenbegriff}

\begin{definition}[Cantorscher Mengenbegriff]
\textbf{Georg Cantor} (1845--1918) formulierte das intuitive Mengenverständnis wie folgt:
\begin{quote}
    \glqq Eine \textbf{Menge} ist eine Zusammenfassung von bestimmten, wohlunterschiedenen Objekten unserer Anschauung oder unseres Denkens (welche die \textbf{Elemente} der Menge genannt werden) zu einem Ganzen.\grqq
\end{quote}
\end{definition}

\begin{bemerkung}[Historische Bedeutung]
Der hier verwendete naive Mengenbegriff lässt insbesondere unendliche Gesamtheiten als mathematische Objekte zu. Obwohl intuitiv eingängig, ist der naive Mengenbegriff nicht axiomatisch exakt; ein uneingeschränkter Umgang führt später zu bekannten Antinomien (z.\,B. der Russellschen Antinomie).
\end{bemerkung}

\subsection*{Beschreibung von Mengen}
Mengen werden üblicherweise auf zwei Arten angegeben:
\begin{enumerate}
    \item \textbf{Aufzählende Form:} $M := \{1, 2, 3\}$ oder $M := \{\text{Hund}, \text{Katze}\}$.
    \item \textbf{Beschreibende Form (charakterisierende Eigenschaft):} 
    \[
    M := \{x \in \Omega \mid H(x)\}, \qquad \text{z.\,B.} \quad M := \{n \in \N \mid 2 \le n \le 5000\}.
    \]
\end{enumerate}

\begin{bemerkung}[Auslassungspunkte]
Die Schreibweise $\{2, 3, \dots, 5000\}$ setzt ein klares Vorwissen über die Bildungsregel voraus. Eine Notation wie $\{1, 3, \dots, 10, 17, 35\}$ ist mathematisch nicht sinnvoll, da die ausgelassenen Elemente unbestimmt sind.
\end{bemerkung}

\begin{beispiel}[Standardmengen]
\begin{itemize}
    \item Leere Menge: $\emptyset = \{\}$ (enthält kein Element)
    \item Natürliche Zahlen: $\N = \{0, 1, 2, 3, \dots\}$
    \item Ganze Zahlen: $\Z = \{\dots, -2, -1, 0, 1, 2, \dots\}$
    \item Rationale Zahlen: $\Q = \left\{ \frac{p}{q} \;\middle|\; p \in \Z, \, q \in \N_{>0} \right\}$
    \item Reelle Zahlen: $\R$
    \item Quadratzahlen: $\{x \in \N \mid \exists y \in \N : y^2 = x\} = \{0, 1, 4, 9, 16, 25, \dots\}$
    \item Abgeschlossenes Intervall: Für $a < b$ sei $[a, b] := \{x \in \R \mid a \le x \le b\}$.
\end{itemize}
\end{beispiel}

% --------------------------------------------------------------------------
% Element- und Teilmengenrelation
% --------------------------------------------------------------------------
\subsubsection*{Element- und Teilmengenrelation}

Für ein Objekt $x$ und eine Menge $M$ gilt entweder $x \in M$ ($x$ ist Element von $M$) oder $x \notin M$, aber niemals beides.

\begin{definition}[Teilmenge und Mengengleichheit]
Seien $M$ und $N$ Mengen.
\begin{enumerate}
    \item $M$ heißt \textbf{Teilmenge} von $N$ (in Zeichen $M \subset N$), falls:
    \[
    \forall x \in M : x \in N.
    \]
    \item $M$ und $N$ heißen \textbf{gleich} (in Zeichen $M = N$), wenn:
    \[
    M \subset N \quad \text{und} \quad N \subset M.
    \]
\end{enumerate}
\end{definition}

\begin{bemerkung}[Wohlunterschiedenheit \& Notation]
\begin{itemize}
    \item Das Symbol $M \subset N$ erlaubt Gleichheit ($M = N$). Für echte Teilmengen schreibt man $M \subsetneq N$.
    \item Aus der Mengengleichheit folgt $\{1, 2, 1\} = \{1, 2\} = \{2, 1\}$. Reihenfolge und Mehrfachnennungen sind irrelevant (anders als bei Tupeln $(1, 2) \ne (2, 1)$ in $\R^2$).
\end{itemize}
\end{bemerkung}

\begin{satz}
Sei $M$ eine Menge und $x$ ein Objekt. Dann gilt:
\[
x \in M \;\Longleftrightarrow\; \{x\} \subset M.
\]
\end{satz}

\begin{proof}
\begin{itemize}
    \item[$\Rightarrow$:] Sei $x \in M$. Die einelementige Menge $\{x\}$ enthält nur $x$, welches nach Annahme in $M$ liegt. Also ist jedes Element aus $\{x\}$ in $M$, somit $\{x\} \subset M$.
    \item[$\Leftarrow$:] Sei $\{x\} \subset M$. Dann sind alle Elemente von $\{x\}$ auch in $M$. Da $x \in \{x\}$, folgt $x \in M$.
\end{itemize}
\end{proof}

% --------------------------------------------------------------------------
% Mengenoperationen
% --------------------------------------------------------------------------
\subsubsection*{Mengenoperationen und Rechenregeln}

\begin{definition}[Operationen auf Mengen]
Seien $M$ und $N$ Mengen.
\begin{enumerate}
    \item \textbf{Vereinigung:} $x \in M \cup N \Longleftrightarrow (x \in M \lor x \in N)$.
    \item \textbf{Schnittmenge:} $x \in M \cap N \Longleftrightarrow (x \in M \land x \in N)$.
    \item \textbf{Komplement / Differenz:} $x \in N \setminus M \Longleftrightarrow (x \in N \land x \notin M)$.
    \item \textbf{Potenzmenge $\Pot(M)$:} Die Menge aller Teilmengen von $M$:
    \[
    A \in \Pot(M) \;\Longleftrightarrow\; A \subset M.
    \]
\end{enumerate}
\end{definition}

\begin{beispiel}[Potenzmengen]
\begin{itemize}
    \item Für $M = \{1, 2\}$ ist $\Pot(M) = \{\emptyset, \{1\}, \{2\}, \{1, 2\}\}$.
    \item Für jede Menge $M$ gilt stets $\emptyset \in \Pot(M)$ und $M \in \Pot(M)$.
    \item Für die leere Menge: $\Pot(\emptyset) = \{\emptyset\} \ne \emptyset$.
    \item Iterierte Potenzmenge: $\Pot(\Pot(\emptyset)) = \{\emptyset, \{\emptyset\}\}$.
\end{itemize}
\end{beispiel}

\begin{satz}[Rechenregeln für Mengenoperationen]
Für beliebige Mengen $M, N, R$ gelten:
\begin{enumerate}
    \item \textbf{Kommutativität:} $M \cup N = N \cup M$ und $M \cap N = N \cap M$.
    \item \textbf{Distributivgesetze:}
    \begin{align*}
        M \cap (N \cup R) &= (M \cap N) \cup (M \cap R), \\
        M \cup (N \cap R) &= (M \cup N) \cap (M \cup R).
    \end{align*}
    \item \textbf{De Morgansche Regeln für Mengen:}
    \begin{align*}
        M \setminus (N \cup R) &= (M \setminus N) \cap (M \setminus R), \\
        M \setminus (N \cap R) &= (M \setminus N) \cup (M \setminus R).
    \end{align*}
\end{enumerate}
\end{satz}

\begin{proof}[Beweis des ersten Distributivgesetzes]
Wir zeigen die Gleichheit über die Äquivalenz der Elementaussagen für ein beliebiges $x$:
\begin{align*}
    x \in M \cap (N \cup R)
    &\overset{\text{Def.\ } \cap}{\Longleftrightarrow} x \in M \land x \in (N \cup R) \\
    &\overset{\text{Def.\ } \cup}{\Longleftrightarrow} x \in M \land (x \in N \lor x \in R) \\
    &\overset{(*)}{\Longleftrightarrow} (x \in M \land x \in N) \lor (x \in M \land x \in R) \\
    &\overset{\text{Def.\ } \cap}{\Longleftrightarrow} x \in (M \cap N) \lor x \in (M \cap R) \\
    &\overset{\text{Def.\ } \cup}{\Longleftrightarrow} x \in (M \cap N) \cup (M \cap R).
\end{align*}
Im Schritt $(*)$ wurde das logische Distributivgesetz für Aussagen aus Kapitel 1 angewendet. Dies zeigt, wie mengentheoretische Gesetze unmittelbar auf die Aussagenlogik zurückgeführt werden.
\end{proof}

\subsubsection*{Übungsblatt 3, Aufgabe 2: Mengenidentitäten}
\notizquelleDatei{handschrift/Uebungsblatt 3_261117_094748.PDF}{1, 6--7}

\noindent\textbf{Aufgabenstellung.}
Beweisen oder widerlegen Sie für alle Mengen $A$, $B$ und $C$:
\begin{enumerate}
    \item $\bigl((A\cap B)\setminus C\bigr)\cup\bigl((B\cap C)\setminus A\bigr)\cup\bigl((A\cap C)\setminus B\bigr)
    = \bigl((A\cap B)\cup(B\cap C)\cup(A\cap C)\bigr)\setminus(A\cap B\cap C)$.
    \item $\bigl((A\cup B)\setminus C\bigr)\cap\bigl((B\cup C)\setminus A\bigr)\cap\bigl((A\cup C)\setminus B\bigr)
    = \bigl((A\cup B)\cap(B\cup C)\cap(A\cup C)\bigr)\setminus(A\cap B\cap C)$.
\end{enumerate}

\begin{proof}[Lösung zu Teil 1 (wahr)]
Für beliebige Mengen $X,Y$ gilt $(X\cap Y)\setminus Z = (X\cap Y)\setminus(X\cap Y\cap Z)$, denn ein Element $x\in X\cap Y$ liegt genau dann nicht in $Z$, wenn es nicht in $X\cap Y\cap Z$ liegt. Wendet man dies auf die drei Terme der linken Seite an (in einer festen Obermenge $U$, z.\,B.\ $U=A\cup B\cup C$, mit Komplement $(\cdot)^c$), so erhält man:
\begin{align*}
    \bigl((A\cap B)\setminus C\bigr)\cup\bigl((B\cap C)\setminus A\bigr)\cup\bigl((A\cap C)\setminus B\bigr)
    &= \bigl((A\cap B)\cap (A\cap B\cap C)^c\bigr) \\
    &\quad \cup \bigl((B\cap C)\cap (A\cap B\cap C)^c\bigr) \\
    &\quad \cup \bigl((A\cap C)\cap (A\cap B\cap C)^c\bigr) \\
    &= \bigl((A\cap B)\cup(B\cap C)\cup(A\cap C)\bigr)\cap (A\cap B\cap C)^c \\
    &= \bigl((A\cap B)\cup(B\cap C)\cup(A\cap C)\bigr)\setminus(A\cap B\cap C).
\end{align*}
Anschaulich beschreibt dies genau diejenigen Elemente, die in mindestens zwei der Mengen $A,B,C$, aber nicht in allen dreien liegen:

\begin{center}
\begin{tikzpicture}[scale=0.9]
    \begin{scope}
        \clip (0,0.65) circle (1.05);
        \fill[primary!25] (-0.65,-0.4) circle (1.05);
    \end{scope}
    \begin{scope}
        \clip (-0.65,-0.4) circle (1.05);
        \fill[primary!25] (0.65,-0.4) circle (1.05);
    \end{scope}
    \begin{scope}
        \clip (0,0.65) circle (1.05);
        \fill[primary!25] (0.65,-0.4) circle (1.05);
    \end{scope}
    \begin{scope}
        \clip (0,0.65) circle (1.05);
        \clip (-0.65,-0.4) circle (1.05);
        \fill[white] (0.65,-0.4) circle (1.05);
    \end{scope}
    \draw[thick, primary] (0,0.65) circle (1.05);
    \draw[thick, primary] (-0.65,-0.4) circle (1.05);
    \draw[thick, primary] (0.65,-0.4) circle (1.05);
    \node at (0,1.95) {$A$};
    \node at (-1.9,-0.7) {$B$};
    \node at (1.9,-0.7) {$C$};
\end{tikzpicture}
\end{center}
\end{proof}

\begin{proof}[Widerlegung zu Teil 2 (falsch)]
Wähle als Gegenbeispiel $A=\{1\}$, $B=\{1\}$ und $C=\emptyset$.
\begin{itemize}
    \item \textbf{Linke Seite:} Es ist $(A\cup B)\setminus C = \{1\}$, $(B\cup C)\setminus A = \emptyset$ und $(A\cup C)\setminus B = \emptyset$. Damit ergibt sich
    \[
        \{1\}\cap\emptyset\cap\emptyset = \emptyset.
    \]
    \item \textbf{Rechte Seite:} Es ist $A\cup B = B\cup C = A\cup C = \{1\}$ sowie $A\cap B\cap C = \emptyset$. Damit ergibt sich
    \[
        (\{1\}\cap\{1\}\cap\{1\})\setminus\emptyset = \{1\}.
    \]
\end{itemize}
Da $\emptyset\ne\{1\}$, ist die Aussage im Allgemeinen falsch. (Tatsächlich ist die linke Seite für beliebige Mengen $A,B,C$ stets leer, da ein Element im Schnitt weder in $C$, noch in $A$, noch in $B$ liegen dürfte, zugleich aber in $A\cup B$.)
\end{proof}

\vspace{0.8em}
\hrule
\vspace{1em}

% Fortsetzung von Kapitel 2. Quelle: kumulative Mitschrift, PDF-S. 27--33.
\subsection{Die Russellsche Antinomie}
\notizquelle{27--28}

Mengen können selbst Elemente anderer Mengen sein. Dabei sind
Element- und Teilmengenrelation sorgfältig zu unterscheiden:
$\{0\}\ne\{\{0\}\}$, und
$\{\{0\},0,\{0,1\}\}$ enthält drei verschiedene Elemente.
Für jede Menge $M$ gilt $M\subset M$; daraus folgt keineswegs $M\in M$.

\begin{satz}[Russellsche Antinomie]
Die uneingeschränkte Forderung, zu jeder Eigenschaft eine Menge aller
Objekte mit dieser Eigenschaft bilden zu dürfen, führt zu einem Widerspruch.
\end{satz}
\begin{proof}
Angenommen, die Zusammenfassung
\[
    R:=\{M\mid M\text{ ist eine Menge und }M\notin M\}
\]
wäre eine Menge. Die definierende Eigenschaft müsste dann auch für $R$
selbst gelten, also
\[
    R\in R\quad\Longleftrightarrow\quad R\notin R.
\]
Ist $R\in R$, so folgt $R\notin R$; ist $R\notin R$, so folgt $R\in R$.
Beide Fälle widersprechen der Annahme.
\end{proof}

\begin{beispiel}[Der Barbier]
Ein Barbier soll genau die Personen rasieren, die sich nicht selbst rasieren.
Gehört er selbst zur betrachteten Personengruppe, so müsste er sich genau
dann selbst rasieren, wenn er sich nicht selbst rasiert. Ein Barbier mit
dieser Eigenschaft kann nicht existieren. Dies ist eine Veranschaulichung
der Antinomie, kein Beweis über tatsächlich existierende Personen.
\end{beispiel}

\begin{bemerkung}[Lügner und Epimenides]
Die Mitschrift nennt auch den Satz eines Kreters: \glqq Alle Kreter
lügen\grqq. Dieser ist ohne weitere Annahmen kein logisches Paradoxon:
Ist er falsch, folgt nur, dass nicht alle Kreter immer lügen.
Der eigentliche Lügner-Satz \glqq Dieser Satz ist falsch\grqq\ richtet
sich dagegen unmittelbar auf seinen eigenen Wahrheitswert.
\end{bemerkung}

\subsection{Auswege}
\notizquelle{28--32}

Die Mitschrift skizziert verschiedene Reaktionen auf solche Antinomien:
Typentheorien trennen verschiedene Stufen von Objekten; eine Beschränkung
auf endliche Mengen wäre für die hier benötigte Analysis zu eng.
Der in der Vorlesung verfolgte Weg ist die \textbf{Axiomatik}.

\begin{definition}[Axiomatischer Ansatz]
Axiome sind festgelegte Grundannahmen einer Theorie. Ein Beweis leitet
weitere Aussagen durch logische Schlussregeln aus diesen Annahmen ab.
Axiome müssen nicht als unmittelbar \glqq offensichtlich wahr\grqq\ gelten;
entscheidend ist, welche Theorie sie beschreiben.
\end{definition}

Für die Mengenlehre verwenden wir den Rahmen von
\textbf{Zermelo--Fraenkel (ZF)}. Mit dem zusätzlichen Auswahlaxiom spricht
man von \textbf{ZFC}. Die Vorlesung gibt hier einen Überblick, keine
vollständige Axiomenliste.

\begin{description}[style=nextline,leftmargin=0pt]
\item[Extensionalitätsaxiom]
Zwei Mengen sind genau dann gleich, wenn sie dieselben Elemente besitzen:
\[
    M=N\quad\Longleftrightarrow\quad
    \forall x:\ (x\in M\Longleftrightarrow x\in N).
\]
Insbesondere gibt es höchstens eine leere Menge. Ihre Existenz folgt aus
den Existenzaxiomen von ZF; Extensionalität allein liefert nur Eindeutigkeit.
\item[Aussonderungsschema]
Ist $M$ bereits eine Menge und $P(x)$ eine zulässige Eigenschaft, dann ist
\[
    \{x\in M\mid P(x)\}
\]
ebenfalls eine Menge. Die Einschränkung auf eine schon vorhandene Menge
ist wesentlich: Eine uneingeschränkte Menge $\{x\mid P(x)\}$ wird nicht
zugelassen.
\item[Weitere Existenzforderungen]
Aus den entsprechenden ZF-Axiomen folgt: Für Mengen $M,N$ sind $M\cup N$
und $\Pot(M)$ Mengen, und es existiert eine unendliche Menge.
Das Auswahlaxiom ist eine zusätzliche Annahme und wird hier nicht
weiter ausgeführt.
\end{description}

\begin{beispiel}[Aussonderung in einer gegebenen Menge]
Der Schnitt ist durch Aussonderung aus $M$ konstruierbar:
\[
    M\cap N=\{x\in M\mid x\in N\}.
\]
Ebenso ist $M\setminus N=\{x\in M\mid x\notin N\}$ eine Menge.
\end{beispiel}

\begin{bemerkung}[Mengen und echte Klassen]
Eine durch eine Eigenschaft beschriebene Gesamtheit von Mengen nennen wir
hier eine \textbf{Klasse}. In ZF ist dies eine abkürzende Sprechweise,
keine zusätzliche Art von Objekten innerhalb der Theorie.
Jede Menge bestimmt eine Klasse, aber nicht jede Klasse ist eine Menge.
Eine Klasse, die keine Menge ist, heißt \textbf{echte Klasse}.
Die Gesamtheit aller Mengen ist eine echte Klasse: Wäre sie eine Menge
$U$, könnte man durch Aussonderung
$\{M\in U\mid M\notin M\}$ bilden und erhielte erneut Russells Widerspruch.
\end{bemerkung}

\subsection{Unvollständigkeit}
\notizquelle{32--33}

Die Mitschrift gibt einen informellen Ausblick auf Gödels
Unvollständigkeitssätze (1931). Die dortige Formulierung
\glqq jede halbwegs interessante Theorie\grqq\ benötigt folgende
Voraussetzungen.

\begin{satz}[Erster Unvollständigkeitssatz, ohne Beweis]
Eine widerspruchsfreie, effektiv axiomatisierbare formale Theorie, die
genügend elementare Arithmetik ausdrücken kann, ist unvollständig:
Es gibt eine Aussage $A$ ihrer Sprache, für die weder $A$ noch $\neg A$
in dieser Theorie beweisbar ist.
\end{satz}

\begin{satz}[Zweiter Unvollständigkeitssatz, ohne Beweis]
Eine widerspruchsfreie, effektiv axiomatisierbare Theorie mit hinreichend
viel Arithmetik kann ihre eigene, in üblicher Weise formalisierte
Widerspruchsfreiheit nicht innerhalb derselben Theorie beweisen.
\end{satz}

\begin{bemerkung}[Einordnung]
\glqq Unvollständig\grqq\ bedeutet nicht, dass eine Aussage weder wahr
noch falsch ist, sondern dass die gewählten Axiome sie nicht entscheiden.
Ein Beweis der Widerspruchsfreiheit in einer stärkeren Theorie ist damit
nicht ausgeschlossen. Die Sätze sind auch kein Nachweis, dass ZF
widersprüchlich ist. Das Buch-Beispiel der Mitschrift dient lediglich
der Anschauung und ersetzt keinen Beweis.
\end{bemerkung}

\noindent{\small Fachliche Präzisierung der Voraussetzungen: P.~Raatikainen,
\emph{Gödel's Incompleteness Theorems}, Stanford Encyclopedia of Philosophy,
\url{https://plato.stanford.edu/entries/goedel-incompleteness/}.}

% Fortsetzung von Kapitel 2. Quelle: kumulative Mitschrift, PDF-S. 34--45.
\subsection{Relationen}
\notizquelle{34--45; S.~42 ist leer}

\subsubsection*{Geordnete Paare und kartesisches Produkt}
\begin{definition}[Kartesisches Produkt]
Für Mengen $X,Y$ heißt
\[
    X\times Y:=\{(x,y)\mid x\in X,\ y\in Y\}
\]
das \textbf{kartesische Produkt}. Ein geordnetes Paar erfüllt
\[
    (x,y)=(x',y')\quad\Longleftrightarrow\quad x=x'\land y=y'.
\]
Anders als bei $\{x,y\}$ kommt es also auf die Reihenfolge an.
\end{definition}

\begin{bemerkung}[Paare als Mengen]
Eine mögliche mengentheoretische Definition ist
$(x,y):=\{\{x\},\{x,y\}\}$.
Für $x\in X$, $y\in Y$ liegt dieses Paar in
$\Pot(\Pot(X\cup Y))$. Daher lässt sich $X\times Y$ durch Aussonderung
aus einer vorhandenen Menge gewinnen.
\end{bemerkung}

\begin{definition}[Relation]
Eine \textbf{Relation zwischen $X$ und $Y$} ist eine Teilmenge
$R\subset X\times Y$. Für $(x,y)\in R$ schreiben wir auch $xRy$.
Ist $X=Y$, heißt $R$ eine \textbf{Relation auf $X$}.
\end{definition}

\subsubsection*{Ordnungsrelationen}
\begin{definition}[Eigenschaften einer Relation]
Eine Relation $R$ auf $X$ heißt
\begin{enumerate}
    \item \textbf{reflexiv}, falls $\forall x\in X:\ xRx$;
    \item \textbf{transitiv}, falls
    $\forall x,y,z\in X:\ (xRy\land yRz)\Rightarrow xRz$;
    \item \textbf{antisymmetrisch}, falls
    $\forall x,y\in X:\ (xRy\land yRx)\Rightarrow x=y$;
    \item \textbf{total}, falls
    $\forall x,y\in X:\ xRy\lor yRx$;
    \item \textbf{symmetrisch}, falls
    $\forall x,y\in X:\ xRy\Rightarrow yRx$.
\end{enumerate}
Eine reflexive, transitive und antisymmetrische Relation heißt
\textbf{Halbordnung}; ist sie zusätzlich total, heißt sie
\textbf{Totalordnung}.
\end{definition}

\begin{beispiel}[Ordnungen und Gegenbeispiele]
\begin{itemize}
    \item $\le$ auf $\N,\Z,\Q,\R$ ist jeweils eine Totalordnung.
    \item $<$ ist transitiv und antisymmetrisch, aber nicht reflexiv.
    In der obigen Definition ist es auch nicht total: Für $x=y$ gilt
    weder $x<y$ noch $y<x$. Es ist eine \emph{strikte} Ordnung.
    \item Auf $\Pot(X)$ ist $\subset$ eine Halbordnung, wobei das Symbol
    wie bisher Gleichheit zulässt. Für zwei verschiedene Elemente
    $a,b\in X$ sind aber $\{a\}$ und $\{b\}$ nicht vergleichbar.
    Bei mindestens zwei Elementen in $X$ ist diese Ordnung also nicht total.
    \item Die echte Teilmengenrelation $\subsetneq$ ist nicht reflexiv.
\end{itemize}
\end{beispiel}

\begin{proof}[Nachweis für die Teilmengenrelation]
Jede Menge $A$ erfüllt $A\subset A$. Aus $A\subset B$ und $B\subset C$
folgt für jedes $x\in A$ zunächst $x\in B$ und dann $x\in C$, also
$A\subset C$. Aus $A\subset B$ und $B\subset A$ folgt nach der
Definition der Mengengleichheit $A=B$.
\end{proof}

\begin{bemerkung}[Symmetrie und Antisymmetrie]
Antisymmetrie ist nicht die Negation von Symmetrie. Die Gleichheitsrelation
ist beispielsweise sowohl symmetrisch als auch antisymmetrisch.
\end{bemerkung}

\subsubsection*{Äquivalenzrelationen und Äquivalenzklassen}
\begin{definition}[Äquivalenzrelation]
Eine reflexive, symmetrische und transitive Relation auf $X$ heißt
\textbf{Äquivalenzrelation}. Dafür verwenden wir häufig das Symbol $\sim$.
Für $x\in X$ heißt
\[
    [x]:=\{y\in X\mid y\sim x\}
\]
die \textbf{Äquivalenzklasse von $x$}. Sie ist eine Teilmenge von $X$.
\end{definition}

\begin{beispiel}
Die Gleichheit auf einer beliebigen Menge ist eine Äquivalenzrelation;
ihre Klassen sind die einelementigen Mengen $\{x\}$.
Auf der Menge der Geraden in $\R^2$ ist \glqq hat dieselbe Richtung\grqq\
eine Äquivalenzrelation. Hier zählen identische Geraden als parallel;
würde man sie ausschließen, fehlte die Reflexivität.
\end{beispiel}

\begin{lemma}[Klassen sind gleich oder disjunkt]
Für eine Äquivalenzrelation auf $X$ und $x,y\in X$ gilt
\[
    [x]=[y]\quad\text{oder}\quad[x]\cap[y]=\emptyset.
\]
Genauer gilt $[x]=[y]\Longleftrightarrow x\sim y$.
\end{lemma}
\begin{proof}
Wegen der Reflexivität gilt $x\in[x]$; keine Klasse ist leer.
Sei nun $z\in[x]\cap[y]$. Dann gelten $z\sim x$ und $z\sim y$,
also nach Symmetrie und Transitivität auch $x\sim y$.
Ist $w\in[x]$, so folgt aus $w\sim x$ und $x\sim y$ auch $w\sim y$,
also $w\in[y]$. Daher ist $[x]\subset[y]$. Analog gilt die umgekehrte
Inklusion. Ein nichtleerer Schnitt erzwingt somit Gleichheit.
Aus $x\sim y$ folgt mit demselben Argument $[x]=[y]$.
Gilt umgekehrt $[x]=[y]$, dann ist $x\in[y]$, also $x\sim y$.
\end{proof}

\begin{definition}[Quotientenmenge]
Die \textbf{Quotientenmenge}
\[
    X/{\sim}:=\{[x]\mid x\in X\}\subset\Pot(X)
\]
ist die Menge aller Äquivalenzklassen. Ihre Elemente sind selbst Mengen.
Die Klassen zerlegen $X$ in paarweise disjunkte, nichtleere Teilmengen.
\end{definition}

\begin{beispiel}[Gerade und ungerade Zahlen]
Auf $\Z$ sei $a\sim b$ genau dann, wenn $a-b$ gerade ist.
Reflexivität folgt aus $a-a=0$, Symmetrie aus
$b-a=-(a-b)$ und Transitivität aus $a-c=(a-b)+(b-c)$.
Es gibt genau zwei Klassen:
\[
    [0]=\{2k\mid k\in\Z\},\qquad
    [1]=\{2k+1\mid k\in\Z\}.
\]
Damit ist $\Z/{\sim}=\{[0],[1]\}$ und $\Z=[0]\cup[1]$.
\end{beispiel}

\begin{beispiel}[Kongruenz modulo $m$]
Für eine feste Zahl $m\in\N_{>0}$ definieren wir auf $\Z$
\[
    a\sim_m b\quad\Longleftrightarrow\quad m\mid(a-b).
\]
Wegen $m\mid0$, $m\mid d\Rightarrow m\mid(-d)$ und
$m\mid d,\ m\mid e\Rightarrow m\mid(d+e)$ ist dies eine
Äquivalenzrelation. Es gilt
\[
    [r]=\{r+km\mid k\in\Z\},\qquad
    \Z/{\sim_m}=\{[0],[1],\ldots,[m-1]\}.
\]
Für $m=3$ sind die Klassen beispielsweise
\begin{align*}
    [0]&=\{\ldots,-6,-3,0,3,6,\ldots\},\\
    [1]&=\{\ldots,-5,-2,1,4,7,\ldots\},\\
    [2]&=\{\ldots,-4,-1,2,5,8,\ldots\}.
\end{align*}
Die untere Zeile der Mitschrift auf PDF-S.~45 verwendet außerdem die
Schreibweise $\Z_m=\{\overline0,\ldots,\overline{m-1}\}$ aus der
Linearen Algebra. Dabei bezeichnen die Symbole Restklassen, keine
gewöhnlichen ganzen Zahlen.
\end{beispiel}

\begin{beispiel}[Planeten-Beispiel von der Tafel]
Unter der Modellannahme, dass jeder betrachtete Planet genau einen Stern
umkreist, sei $P\sim Q$ genau dann, wenn beide denselben Stern umkreisen.
Die Relation ist reflexiv und symmetrisch. Umkreisen $P,Q$ denselben
Stern $S_1$ und $Q,R$ denselben Stern $S_2$, so folgt aus der Eindeutigkeit
für $Q$, dass $S_1=S_2$. Daher gilt $P\sim R$.
Die Klassen sind die Gruppen von Planeten um denselben Stern.
\end{beispiel}

% Fortsetzung von Kapitel 2. Quelle: kumulative Mitschrift, PDF-S. 46--54.
\subsection{Abbildungen}
\notizquelle{46--54}

\subsubsection*{Definition und Wohldefiniertheit}
\begin{definition}[Abbildung und Graph]
Eine \textbf{Abbildung} $f:X\to Y$ ordnet jedem $x\in X$ genau ein
Element $f(x)\in Y$ zu. Wir schreiben $x\mapsto f(x)$.
$X$ heißt \textbf{Definitionsbereich}, $Y$ heißt \textbf{Zielbereich}.
Formal wird die Zuordnung durch eine Relation $G_f\subset X\times Y$
mit der Eigenschaft
\[
    \forall x\in X\;\exists!y\in Y:\ (x,y)\in G_f
\]
beschrieben. $G_f=\{(x,f(x))\mid x\in X\}$ ist der \textbf{Graph} von $f$.
Definitions- und Zielbereich gehören zur Angabe der Abbildung.
\end{definition}

\begin{beispiel}[Einfache Zuordnungen]
\begin{itemize}
    \item Die Identität $\operatorname{id}_X:X\to X$, $x\mapsto x$,
    hat den Graphen $\{(x,x)\mid x\in X\}$.
    \item $f:\N\to\N$, $n\mapsto n+1$, ist eine Abbildung.
    \item Die Vorschrift $\N\to\N$, $n\mapsto n-1$, ist wegen
    $0\mapsto-1\notin\N$ keine Abbildung in den angegebenen Zielbereich.
    Auf $\Z\to\Z$ ist dieselbe Vorschrift wohldefiniert.
\end{itemize}
Hier und im gesamten Skript gilt $0\in\N$.
\end{beispiel}

\begin{beispiel}[Abhängigkeit von der Darstellung]
Die Vorschrift $\Q\to\Z$, $p/q\mapsto p$, ist ohne Einschränkung
der Bruchdarstellung nicht wohldefiniert: $1/2=2/4$, aber $1\ne2$.
Verlangt man dagegen eine vollständig gekürzte Darstellung mit
$q>0$, ist der Zähler eindeutig. Für $0$ ist die gekürzte Darstellung
$0/1$. Damit erhält man eine wohldefinierte Abbildung nach $\Z$.
Die Positivitätsforderung betrifft den \emph{Nenner}; der Zähler
kann negativ oder null sein.
\end{beispiel}

\begin{beispiel}[Repräsentanten von Restklassen]
Für $m\ge1$ ist die Vorschrift
\[
    \Z/{\sim_m}\longrightarrow\Z,\qquad [k]\longmapsto k
\]
nicht wohldefiniert: $[0]=[m]$, aber $0\ne m$.
Dagegen ist die natürliche Projektion
$\Z\to\Z/{\sim_m}$, $k\mapsto[k]$, wohldefiniert.
Auch die Wahl des \emph{eindeutigen} Repräsentanten in
$\{0,\ldots,m-1\}$ definiert eine Abbildung aus der Quotientenmenge;
dies ist eine zusätzliche, ausdrücklich festgelegte Regel.
\end{beispiel}

\begin{beispiel}[Stückweise Definition und Collatz-Folge]
Die Vorschrift
\[
    c:\N\to\N,\qquad
    c(n)=\begin{cases}
        n/2,&n\text{ gerade},\\
        3n+1,&n\text{ ungerade}
    \end{cases}
\]
ist wohldefiniert, weil genau einer der Fälle zutrifft und das Ergebnis
jeweils in $\N$ liegt. Iteration liefert für den Startwert $7$:
\[
\begin{aligned}
    7,22,11,34,17,52,26,13,40,20,\\
    10,5,16,8,4,2,1,\ldots
\end{aligned}
\]
Die in der Vorlesung vorgestellte Collatz-Vermutung besagt, dass die
Iteration bei jedem \emph{positiven} Startwert irgendwann $1$ erreicht.
Danach wiederholt sich $1,4,2,1,\ldots$; die Folge endet also nicht
von selbst. Der Startwert $0$ bleibt bei $0$ und ist ausdrücklich
aus der Vermutung ausgeschlossen. Wohldefiniertheit der Vorschrift
ist von dieser Aussage über ihr langfristiges Verhalten zu unterscheiden.
\end{beispiel}

\subsubsection*{Komposition und Einschränkung}
\begin{definition}[Komposition]
Für Abbildungen $f:X\to Y$ und $g:Y\to Z$ ist die
\textbf{Komposition} $g\circ f:X\to Z$ gegeben durch
\[
    (g\circ f)(x):=g(f(x)).
\]
Man wendet zuerst $f$ und dann $g$ an.
\end{definition}

\begin{beispiel}[Die Reihenfolge ist wesentlich]
Für $f,g:\R\to\R$, $f(x)=x+1$ und $g(x)=x^2$, gilt
\[
    (g\circ f)(x)=(x+1)^2,
    \qquad (f\circ g)(x)=x^2+1.
\]
Bei $x=1$ ergeben sich $4$ bzw.\ $2$, also $g\circ f\ne f\circ g$.
Im Allgemeinen müssen nicht einmal beide Kompositionen definiert sein.
\end{beispiel}

\begin{definition}[Einschränkung]
Für $A\subset X$ ist die Inklusion $\iota:A\to X$, $x\mapsto x$,
eine Abbildung. Die \textbf{Einschränkung} von $f:X\to Y$ auf $A$ ist
\[
    f|_A:=f\circ\iota:A\to Y,\qquad x\mapsto f(x).
\]
\end{definition}

\subsubsection*{Graphen, Bilder und Urbilder}
\begin{beispiel}[Zwei Funktionsgraphen]
Die Betragsfunktion und die Sprungfunktion der Mitschrift sind
\[
    |x|=\begin{cases}x,&x\ge0,\\-x,&x<0,\end{cases}
    \qquad
    h(x)=\begin{cases}0,&x\le0,\\1,&x>0.\end{cases}
\]
Beide sind Abbildungen $\R\to\R$; insbesondere ist $h(0)=0$.
\begin{center}
\begin{tikzpicture}[scale=0.72]
    \draw[->] (-2.5,0)--(2.7,0) node[right] {$x$};
    \draw[->] (0,-0.5)--(0,2.8) node[above] {$y$};
    \draw[thick,primary] (-2.2,2.2)--(0,0)--(2.2,2.2);
    \node at (0,-1) {$y=|x|$};
\end{tikzpicture}
\hspace{1.2cm}
\begin{tikzpicture}[scale=0.72]
    \draw[->] (-2.5,0)--(2.7,0) node[right] {$x$};
    \draw[->] (0,-0.5)--(0,2.8) node[above] {$y$};
    \draw[very thick,primary] (-2.2,0)--(0,0);
    \draw[very thick,primary] (0,1)--(2.2,1);
    \fill[primary] (0,0) circle (2.8pt);
    \draw[primary,fill=white,thick] (0,1) circle (2.8pt);
    \node[left] at (0,1) {$1$};
    \node at (0,-1) {$y=h(x)$};
\end{tikzpicture}
\end{center}
\end{beispiel}

\begin{definition}[Bild und Urbild]
Sei $f:X\to Y$ eine Abbildung.
\begin{enumerate}
    \item Für $A\subset X$ heißt
    \[
        f(A):=\{f(x)\mid x\in A\}
        =\{y\in Y\mid\exists x\in A:\ f(x)=y\}
    \]
    das \textbf{Bild von $A$}. Das \textbf{Bild der Abbildung} ist
    $\operatorname{im}(f):=f(X)\subset Y$.
    \item Für $B\subset Y$ heißt
    \[
        f^{-1}(B):=\{x\in X\mid f(x)\in B\}
    \]
    das \textbf{Urbild von $B$}.
    Für ein einzelnes $y\in Y$ schreibt man auch
    $f^{-1}(y):=f^{-1}(\{y\})$ und nennt dies die \textbf{Faser über $y$}.
\end{enumerate}
\end{definition}

\begin{bemerkung}[Urbild ist keine Umkehrfunktion]
Urbilder sind für jede Abbildung definiert; sie sind Mengen und dürfen
leer sein oder mehrere Elemente enthalten. Die Schreibweise $f^{-1}(B)$
setzt keine Umkehrfunktion voraus. Außerdem ist $f(x)$ ein Element,
während $f(A)$ eine Menge ist. Der Zielbereich $Y$ kann größer als $f(X)$ sein.
\end{bemerkung}

\begin{beispiel}
Für $q:\Z\to\N$, $q(n)=n^2$, sind
$q(\{-2,0,2\})=\{0,4\}$,
$q^{-1}(\{4\})=\{-2,2\}$ und $q^{-1}(\{5\})=\emptyset$.
\end{beispiel}

\subsubsection*{Injektivität, Surjektivität und Bijektivität}
\begin{definition}
Eine Abbildung $f:X\to Y$ heißt
\begin{enumerate}
    \item \textbf{injektiv}, falls
    $\forall x,x'\in X:\ f(x)=f(x')\Rightarrow x=x'$;
    \item \textbf{surjektiv}, falls
    $\forall y\in Y\;\exists x\in X:\ f(x)=y$;
    \item \textbf{bijektiv}, falls sie injektiv und surjektiv ist.
\end{enumerate}
In Fasern formuliert: Für jedes $y\in Y$ enthält $f^{-1}(\{y\})$
bei Injektivität höchstens ein, bei Surjektivität mindestens ein und
bei Bijektivität genau ein Element.
\end{definition}

\begin{beispiel}[Beispiele aus der Mitschrift]
\begin{enumerate}
    \item $\{1,2,3\}\to\{2,4,6,8\}$, $n\mapsto2n$, ist injektiv,
    aber wegen der fehlenden $8$ nicht surjektiv.
    \item $\N\to\{35\}$, $n\mapsto35$, ist surjektiv, aber nicht
    injektiv, etwa wegen $f(0)=f(1)$.
    \item $\Z\to\N$, $n\mapsto n^2$, ist weder injektiv
    ($2^2=(-2)^2$) noch surjektiv ($5$ ist kein Quadrat).
    \item $\Z\to\N$, $n\mapsto|n|$, ist surjektiv, da $|k|=k$
    für jedes $k\in\N$, aber wegen $|-2|=|2|$ nicht injektiv.
    \item $\Z\to\Z$, $n\mapsto-n$, ist bijektiv:
    Aus $-a=-b$ folgt $a=b$, und $-(-y)=y$ liefert zu jedem $y$ ein Urbild.
\end{enumerate}
\end{beispiel}

\begin{satz}[Zielbereich auf das Bild einschränken]
Für jede Abbildung $f:X\to Y$ ist
$\widetilde f:X\to f(X)$, $x\mapsto f(x)$, surjektiv.
\end{satz}
\begin{proof}
Jedes $y\in f(X)$ hat nach Definition die Form $y=f(x)$ für ein $x\in X$.
\end{proof}

\begin{bemerkung}[Wohldefiniertheit zuerst]
Wohldefiniertheit verlangt, dass \emph{jedes Eingabeelement genau einen
Ausgabewert} im Zielbereich erhält. Injektivität verlangt dagegen,
dass \emph{verschiedene Eingaben verschiedene Ausgaben} haben.
Eine konstante Abbildung ist wohldefiniert und im Allgemeinen nicht injektiv.
Surjektivität hängt vom angegebenen Zielbereich ab: $n\mapsto2n$
ist von $\N$ nach $\N$ nicht surjektiv, von $\N$ nach
$\{2n\mid n\in\N\}$ dagegen schon.
\end{bemerkung}

\notizquelleNeu{2--7}

\begin{bemerkung}[Logische Verneinung der Injektivität]
Aus der Definition ergibt sich durch Verneinung und Kontraposition:
\begin{align*}
    f\text{ nicht injektiv}
    &\iff \exists x_1,x_2\in X:\ \bigl(f(x_1)=f(x_2)\land x_1\ne x_2\bigr),\\
    f\text{ injektiv}
    &\iff \neg\bigl(\exists x_1,x_2\in X:\ (f(x_1)=f(x_2)\land x_1\ne x_2)\bigr)\\
    &\iff \forall x_1,x_2\in X:\ \bigl(f(x_1)\ne f(x_2)\lor x_1=x_2\bigr)\\
    &\iff \forall x_1,x_2\in X:\ \bigl(f(x_1)=f(x_2)\implies x_1=x_2\bigr),
\end{align*}
wobei im letzten Schritt $(A\implies B)\equiv(\neg A\lor B)$ eingeht.
In Pfeildiagrammen zwischen endlichen Mengen lässt sich dies unmittelbar
ablesen:
\begin{center}
\begin{tikzpicture}[scale=0.82]
    % Keine Abbildung
    \draw (0,0) ellipse (0.55 and 1.15);
    \draw (2,0) ellipse (0.55 and 1.15);
    \fill (0,0.6) circle (1.5pt) (0,0) circle (1.5pt) (0,-0.6) circle (1.5pt);
    \fill (2,0.75) circle (1.5pt) (2,0.25) circle (1.5pt) (2,-0.25) circle (1.5pt) (2,-0.75) circle (1.5pt);
    \draw[->] (0.08,0.6)--(1.92,0.75);
    \draw[->] (0.08,0)--(1.92,0.25);
    \draw[->] (0.08,0)--(1.92,-0.25);
    \draw[->] (0.08,-0.6)--(1.92,-0.75);
    \node[align=center,font=\small] at (1,-1.55) {keine Abbildung!};

    % Injektiv, nicht surjektiv
    \begin{scope}[xshift=4.2cm]
    \draw (0,0) ellipse (0.55 and 1.15);
    \draw (2,0) ellipse (0.55 and 1.15);
    \fill (0,0.5) circle (1.5pt) (0,0) circle (1.5pt) (0,-0.5) circle (1.5pt);
    \fill (2,0.75) circle (1.5pt) (2,0.25) circle (1.5pt) (2,-0.25) circle (1.5pt) (2,-0.75) circle (1.5pt);
    \draw[->] (0.08,0.5)--(1.92,0.75);
    \draw[->] (0.08,0)--(1.92,0.25);
    \draw[->] (0.08,-0.5)--(1.92,-0.25);
    \node[align=center,font=\small] at (1,-1.55) {injektiv,\\nicht surjektiv};
    \end{scope}

    % Surjektiv, nicht injektiv
    \begin{scope}[xshift=8.4cm]
    \draw (0,0) ellipse (0.55 and 1.15);
    \draw (2,0) ellipse (0.55 and 1.15);
    \fill (0,0.7) circle (1.5pt) (0,0.2) circle (1.5pt) (0,-0.3) circle (1.5pt) (0,-0.75) circle (1.5pt);
    \fill (2,0.5) circle (1.5pt) (2,0) circle (1.5pt) (2,-0.5) circle (1.5pt);
    \draw[->] (0.08,0.7)--(1.92,0.5);
    \draw[->] (0.08,0.2)--(1.92,0.5);
    \draw[->] (0.08,-0.3)--(1.92,0);
    \draw[->] (0.08,-0.75)--(1.92,-0.5);
    \node[align=center,font=\small] at (1,-1.55) {surjektiv,\\nicht injektiv};
    \end{scope}
\end{tikzpicture}
\end{center}
Ebenso ist $f:\R\to\R_{\ge0}$, $x\mapsto x^2$, surjektiv, aber nicht injektiv,
etwa wegen $f^{-1}(\{2\})=\{\sqrt{2},-\sqrt{2}\}$.
\end{bemerkung}

\begin{lemma}[Komposition injektiver und surjektiver Abbildungen]
Seien $f:X\to Y$ und $g:Y\to Z$ Abbildungen zwischen Mengen.
\begin{enumerate}
    \item Sind $f$ und $g$ injektiv, dann ist auch $g\circ f:X\to Z$ injektiv.
    \item Sind $f$ und $g$ surjektiv, dann ist auch $g\circ f:X\to Z$ surjektiv.
    \item Sind $f$ und $g$ bijektiv, dann ist auch $g\circ f:X\to Z$ bijektiv.
\end{enumerate}
\end{lemma}
\begin{proof}
\begin{enumerate}
    \item Seien $x_1,x_2\in X$ mit $(g\circ f)(x_1)=(g\circ f)(x_2)$, also
    $g(f(x_1))=g(f(x_2))$. Da $g$ injektiv ist, folgt $f(x_1)=f(x_2)$.
    Da $f$ injektiv ist, folgt daraus $x_1=x_2$.
    \item Sei $z\in Z$. Da $g$ surjektiv ist, existiert ein $y\in Y$ mit
    $g(y)=z$. Da $f$ surjektiv ist, existiert ein $x\in X$ mit $f(x)=y$.
    Somit gilt $(g\circ f)(x)=g(f(x))=g(y)=z$.
    \item Folgt unmittelbar aus (1) und (2). \qedhere
\end{enumerate}
\end{proof}

\subsubsection*{Umkehrabbildung}
\begin{definition}[Umkehrabbildung]
Sei $f:X\to Y$ eine bijektive Abbildung. Dann ist für jedes Element
$y\in Y$ das Urbild $f^{-1}(\{y\})=\{x\in X\mid f(x)=y\}$ eine
einelementige Menge. Wir definieren die \textbf{Umkehrabbildung}
$f^{-1}:Y\to X$, indem jedes $y\in Y$ auf das eindeutige Element in
$f^{-1}(\{y\})$ geschickt wird; mit anderen Worten: $f^{-1}(y)\in X$ ist
das eindeutige Element in $X$ mit $f(f^{-1}(y))=y$.
\end{definition}

\begin{bemerkung}[Urbild vs.\ Umkehrabbildung]
Die Schreibweise $f^{-1}(y)$ bezeichnet damit zwei verschiedene Dinge:
Das Urbild $f^{-1}(\{y\})$ (oft kurz $f^{-1}(y)$ geschrieben) ist für
\emph{jede} Abbildung als Teilmenge von $X$ definiert, die Umkehrabbildung
$f^{-1}:Y\to X$ hingegen nur für \emph{bijektive} Abbildungen $f$.
\end{bemerkung}

\begin{beispiel}
Die Abbildung $f:\R_{\ge0}\to\R_{\ge0}$, $x\mapsto x^2$, ist bijektiv.
Ihre Umkehrabbildung ist $f^{-1}:\R_{\ge0}\to\R_{\ge0}$, $x\mapsto\sqrt{x}$.
\end{beispiel}

\begin{satz}[Charakterisierung der Umkehrabbildung]
\begin{enumerate}
    \item Ist $f:X\to Y$ bijektiv, dann erfüllt die Umkehrabbildung
    $f^{-1}:Y\to X$ die Gleichungen
    \[
        f\circ f^{-1}=\operatorname{id}_Y
        \qquad\text{und}\qquad
        f^{-1}\circ f=\operatorname{id}_X.
    \]
    \item Sind $f:X\to Y$ und $g:Y\to X$ Abbildungen mit
    $f\circ g=\operatorname{id}_Y$ und $g\circ f=\operatorname{id}_X$,
    dann ist $f$ bijektiv und es gilt $g=f^{-1}$.
\end{enumerate}
\end{satz}
\begin{proof}
\begin{enumerate}
    \item Für jedes $y\in Y$ gilt $f(f^{-1}(y))=y$ nach Definition von $f^{-1}$,
    also $f\circ f^{-1}=\operatorname{id}_Y$. Für jedes $x\in X$ liegt $x$
    im Urbild $f^{-1}(\{f(x)\})$, und da dieses einelementig ist, gilt
    $f^{-1}(f(x))=x$, also $f^{-1}\circ f=\operatorname{id}_X$.
    \item Aus $g\circ f=\operatorname{id}_X$ folgt die Injektivität von $f$:
    Gilt $f(x_1)=f(x_2)$, so erhält man durch Anwenden von $g$ sofort
    $x_1=g(f(x_1))=g(f(x_2))=x_2$.
    Aus $f\circ g=\operatorname{id}_Y$ folgt die Surjektivität von $f$:
    Zu $y\in Y$ liefert $x:=g(y)\in X$ gerade $f(x)=f(g(y))=y$.
    Also ist $f$ bijektiv. Schließlich gilt für jedes $y\in Y$:
    $g(y)=f^{-1}(f(g(y)))=f^{-1}(y)$, also $g=f^{-1}$. \qedhere
\end{enumerate}
\end{proof}

% Fortsetzung von Kapitel 2 (§ 7). Quelle: vorlesungsmitschrift_gdm1PDF_261114_161111.PDF, PDF-S. 7--13.
\subsection{Das Auswahlaxiom}
\notizquelleNeu{7--13}

\begin{definition}[Auswahlaxiom (AC)]
Das \textbf{Auswahlaxiom} (\emph{Axiom of Choice}, kurz \textbf{AC}) besagt:
Zu jeder nichtleeren Menge $X$ von nichtleeren Mengen existiert eine
\textbf{Auswahlfunktion}, das heißt eine Abbildung
\[
    f:X\longrightarrow\bigcup_{x\in X}x,
    \qquad\text{sodass für alle }x\in X\text{ gilt: }f(x)\in x.
\]
\end{definition}

\begin{beispiel}[Auswahlfunktionen]
\begin{enumerate}
    \item Für $X=\{\{1,2\},\{1\},\{3\},\{1,2,4\}\}$ ist
    $f:X\to\{1,2,3,4\}$ mit
    \[
        \{1,2\}\mapsto1,\quad
        \{1\}\mapsto1,\quad
        \{3\}\mapsto3,\quad
        \{1,2,4\}\mapsto4
    \]
    eine Auswahlfunktion für $X$.
    \item Für die Menge der nichtleeren Teilmengen der natürlichen Zahlen
    $X=\Pot(\N)\setminus\{\emptyset\}$ liefert
    \[
        f:\Pot(\N)\setminus\{\emptyset\}\longrightarrow\N,\qquad
        Z\longmapsto\min(Z)
    \]
    (die Zuordnung zum Minimum der Elemente von $Z$) eine explizite
    Auswahlfunktion. (In der Mitschrift kurz als $\Pot(\N)$ notiert;
    für die leere Menge existiert kein Minimum.)
    \item Für $X=\Pot(\R)\setminus\{\emptyset\}$ ist es nicht möglich,
    explizit eine Auswahlfunktion hinzuschreiben; das Auswahlaxiom
    garantiert jedoch ihre Existenz.
    \item \emph{Schuhe und Socken (Russell):} Ist $X$ eine Menge von
    Paaren von Schuhen, so lässt sich eine Auswahlfunktion explizit
    angeben, indem aus jedem Paar der linke Schuh gewählt wird.
    Ist $X$ dagegen eine Menge ununterscheidbarer Paare von Socken,
    lässt sich keine solche Regel angeben -- hier benötigt man AC.
\end{enumerate}
\end{beispiel}

\begin{bemerkung}[ZF, ZFC und Unabhängigkeit]
Das Auswahlaxiom wird manchmal zum Zermelo-Fraenkelschen Axiomensystem
dazugezählt und manchmal nicht: Man unterscheidet zwischen $\mathrm{ZF}$
und $\mathrm{ZFC}$ ($\mathrm{ZF}+\mathrm{AC}$).
Nach Resultaten von Gödel und Cohen ist AC unabhängig von $\mathrm{ZF}$:
Angenommen, $\mathrm{ZF}$ ist widerspruchsfrei, dann ist sowohl
$\mathrm{ZFC}$ als auch $\mathrm{ZF}+\neg\mathrm{AC}$ widerspruchsfrei.
\end{bemerkung}

\subsubsection*{Ordnungen und das Lemma von Zorn}
\begin{definition}[Halbordnung und Totalordnung]
Eine Relation $\le$ auf einer Menge $X$ heißt \textbf{Halbordnung}, wenn sie
\begin{enumerate}
    \item \textbf{reflexiv} ist: $\forall x\in X:\ x\le x$;
    \item \textbf{transitiv} ist: $\forall x,y,z\in X:\ (x\le y\land y\le z)\implies x\le z$;
    \item \textbf{antisymmetrisch} ist: $\forall x,y\in X:\ (x\le y\land y\le x)\implies x=y$.
\end{enumerate}
Eine Halbordnung heißt genau dann \textbf{Totalordnung}, wenn sie total ist,
das heißt $\forall x,y\in X:\ (x\le y\lor y\le x)$.
\end{definition}

\begin{beispiel}
Für jede Menge $X$ ist die Teilmengenrelation $\subset$ auf der
Potenzmenge $\Pot(X)$ eine Halbordnung (im Allgemeinen keine Totalordnung).
\end{beispiel}

\begin{definition}[Extrema, minimale/maximale Elemente und Schranken]
Sei $\le$ eine Halbordnung auf einer Menge $X$ und $Y\subset X$ eine Teilmenge.
\begin{itemize}
    \item Ein \textbf{Minimum} (bzw.\ \textbf{Maximum}) von $Y$ ist ein Element
    $y\in Y$, sodass für alle $y'\in Y$ gilt: $y\le y'$ (bzw.\ $y'\le y$).
    \item Ein \textbf{minimales Element} (bzw.\ \textbf{maximales Element}) von $Y$
    ist ein Element $y\in Y$, sodass es kein $y'\in Y$ gibt mit $y'\ne y$ und
    $y'\le y$ (bzw.\ $y\le y'$).
    \item Eine \textbf{untere} (bzw.\ \textbf{obere}) \textbf{Schranke} von $Y$
    ist ein Element $x\in X$, sodass für alle $y\in Y$ gilt:
    $x\le y$ (bzw.\ $y\le x$).
\end{itemize}
\end{definition}

\begin{bemerkung}
Ist $y\in Y$ ein Maximum von $Y$, dann ist $y$ sowohl ein maximales Element
von $Y$ als auch eine obere Schranke von $Y$.
Wenn ein Maximum (bzw.\ Minimum) von $Y$ existiert, ist es eindeutig:
Sind nämlich $y_1,y_2\in Y$ Maxima, so gilt $y_1\le y_2$ und $y_2\le y_1$,
also $y_1=y_2$ wegen der Antisymmetrie.
\end{bemerkung}

\begin{beispiel}[Obere Schranken und maximale Elemente]
\begin{enumerate}
    \item Auf $X=\Q$ mit der üblichen Ordnung $\le$ hat das offene Intervall
    $Y=(0,1)=\{x\in\Q\mid 0<x<1\}$ die Zahl $1$ als obere Schranke, aber
    $1\notin Y$ ist kein Maximum von $Y$. Auch $2$ ist eine obere Schranke;
    obere Schranken sind im Allgemeinen nicht eindeutig.
    \item Auf $Y=\{\{1,2\},\{1\},\{3\},\{1,2,4\}\}\subset\Pot(\N)$ mit der
    Teilmengenrelation $\subset$ gelten
    $\{1\}\subset\{1,2\}\subset\{1,2,4\}$ und $\{1\}\subset\{1,2,4\}$.
    Daher sind sowohl $\{1,2,4\}$ als auch $\{3\}$ maximale Elemente von $Y$
    (keines ist echt in einer anderen Menge aus $Y$ enthalten). Ein Maximum
    existiert nicht, da $\{3\}$ und $\{1,2,4\}$ unvergleichbar sind.
    Die Teilmenge $\{\{1\},\{1,2\},\{1,2,4\}\}\subset Y$ ist eine
    totalgeordnete Teilmenge von $Y$.
\end{enumerate}
\end{beispiel}

\begin{satz}[Lemma von Zorn]
Eine halbgeordnete Menge $(X,\le)$, in der jede totalgeordnete Teilmenge
(das heißt je zwei Elemente in der Teilmenge sind vergleichbar) eine
obere Schranke in $X$ besitzt, besitzt ein maximales Element.
\end{satz}

\begin{bemerkung}[Äquivalenz zu AC]
Im Rahmen von $\mathrm{ZF}$ sind das Auswahlaxiom und das Lemma von Zorn
äquivalent: $\mathrm{AC}\iff\text{Zorn}$ (ohne Beweis).
(Da auch die leere Teilmenge $\emptyset\subset X$ totalgeordnet ist,
impliziert die Voraussetzung des Lemmas von Zorn insbesondere $X\ne\emptyset$.)
\end{bemerkung}

% Fortsetzung von Kapitel 2 (§ 8). Quelle: vorlesungsmitschrift_gdm1PDF_261114_161111.PDF, PDF-S. 14--23.
\subsection{Mächtigkeit von Mengen}
\vorlesungsdatum{13.\ November 2026}{vorlesungsmitschrift_gdm1PDF_261114_161111.PDF}
\notizquelleNeu{14--23}

\begin{definition}[Vergleich von Mächtigkeiten]
Seien $M$ und $N$ Mengen. Dann sagen wir:
\begin{enumerate}
    \item Die Mächtigkeit von $M$ ist kleiner oder gleich der Mächtigkeit
    von $N$ (geschrieben $|M|\le|N|$), wenn es eine injektive Abbildung
    $f:M\to N$ gibt.
    \item $M$ und $N$ sind \textbf{gleichmächtig} (geschrieben $|M|=|N|$),
    wenn es eine bijektive Abbildung $f:M\to N$ gibt.
\end{enumerate}
Man schreibt $|M|<|N|$ genau dann, wenn $|M|\le|N|$ und $|M|\ne|N|$ gilt.
\end{definition}

\begin{bemerkung}[Endliche und unendliche Mengen]
Das Symbol $|M|$ hat für allgemeine Mengen keine eigenständige Bedeutung;
definiert sind zunächst nur die Relationen $|M|\le|N|$, $|M|=|N|$ und
$|M|<|N|$.
Ist $M$ endlich, so bezeichnet $|M|$ die Anzahl der Elemente von $M$,
etwa $|\{1,2,4,7\}|=4$ oder $|\{\{1\},\{1,3\}\}|=2$.
Eine Menge $M$ ist genau dann \textbf{unendlich}, wenn sie gleichmächtig
zu einer echten Teilmenge von sich selbst ist.
\end{bemerkung}

\begin{beispiel}[Erste Gleichmächtigkeiten und Vergleiche]
\begin{enumerate}
    \item $\N\to\N\setminus\{0\}$, $n\mapsto n+1$, ist eine Bijektion,
    also $|\N|=|\N\setminus\{0\}|$.
    \item $\N\to2\N=\{\text{gerade natürliche Zahlen}\}$, $n\mapsto 2n$,
    ist eine Bijektion, also $|\N|=|2\N|$.
    \item Durch Inklusionen gilt $|\N|\le|\Z|\le|\Q|\le|\R|\le|\C|$.
    \item Für jede Menge $M$ gilt $|M|\le|\Pot(M)|$, denn die Abbildung
    $M\to\Pot(M)$, $x\mapsto\{x\}$, ist injektiv.
\end{enumerate}
\end{beispiel}

\begin{definition}[Abzählbar und überabzählbar]
Eine Menge $M$ heißt \textbf{abzählbar} (oder \textbf{abzählbar unendlich}),
wenn $|M|=|\N|$ gilt.
Sie heißt \textbf{überabzählbar}, wenn $|\N|<|M|$ gilt.
\end{definition}

\begin{satz}[Abzählbarkeit von $\Z$]
Es gilt $|\N|=|\Z|$.
\end{satz}
\begin{proof}
Die Abbildung
\[
    f:\N\longrightarrow\Z,\qquad
    n\longmapsto\begin{cases}
        \dfrac{n}{2}, & n\text{ gerade},\\[0.6em]
        -\dfrac{n+1}{2}, & n\text{ ungerade},
    \end{cases}
\]
ist eine Bijektion: Die ersten Werte sind
$0\mapsto0$, $1\mapsto-1$, $2\mapsto1$, $3\mapsto-2$, $4\mapsto2$,
$5\mapsto-3$, $6\mapsto3$, \dots; jedes $k\ge0$ hat genau das Urbild $2k$,
jedes $k<0$ genau das Urbild $-2k-1$.
\end{proof}

\begin{bemerkung}[Endliche vs.\ unendliche Produkte]
Sind $M$ und $N$ endliche Mengen mit $|M|=m\ge2$ und $|N|=n\ge2$,
so hat $M\times N$ genau $m\cdot n$ Elemente; wegen $m\cdot n>m$ und
$m\cdot n>n$ gilt dann $|M\times N|>|M|$ und $|M\times N|>|N|$.
Bei unendlichen Mengen verhält sich das kartesische Produkt anders.
\end{bemerkung}

\begin{satz}[Abzählbarkeit von $\N\times\N$]
Es gilt $|\N\times\N|=|\N|$.
\end{satz}
\begin{proof}
Wegen $|\N|=|\N\setminus\{0\}|$ genügt es zu zeigen, dass
$|\N\setminus\{0\}|=|(\N\setminus\{0\})\times(\N\setminus\{0\})|$.
Dies gilt, da die Abbildung
\[
    f:(\N\setminus\{0\})\times(\N\setminus\{0\})\longrightarrow\N\setminus\{0\},
    \qquad (n,m)\longmapsto 2^{n-1}\cdot(2m-1)
\]
bijektiv ist (jede positive natürliche Zahl lässt sich eindeutig als
Produkt einer Zweierpotenz $2^{n-1}$ mit einer ungeraden positiven Zahl
$2m-1$ schreiben).
\end{proof}

\begin{satz}[Injektionen vs.\ Surjektionen bei Mächtigkeiten]
Seien $M$ und $N$ nichtleere Mengen. Dann gilt genau dann $|M|\le|N|$,
wenn es eine surjektive Abbildung $N\to M$ gibt.
\end{satz}
\begin{proof}
\emph{($\Rightarrow$):} Es gelte $|M|\le|N|$, das heißt es gibt eine
injektive Abbildung $f:M\to N$. Wir fassen $f$ als bijektive Abbildung
$\widetilde f:M\to f(M)$ auf das Bild $f(M)=\{f(x)\mid x\in M\}$ auf.
Diese besitzt eine bijektive Umkehrabbildung $\widetilde f^{-1}:f(M)\to M$.
Wähle ein beliebiges $x_0\in M$ (dies ist möglich, da $M\ne\emptyset$)
und definiere
\[
    g:N\longrightarrow M,\qquad
    y\longmapsto\begin{cases}
        \widetilde f^{-1}(y), & y\in f(M),\\
        x_0, & y\notin f(M).
    \end{cases}
\]
Dann ist $g$ surjektiv, denn für jedes $x\in M$ gilt $f(x)\in f(M)$ und
$g(f(x))=\widetilde f^{-1}(f(x))=x$.

\medskip
\noindent\emph{($\Leftarrow$):} Es existiere umgekehrt eine surjektive
Abbildung $g:N\to M$. Zu zeigen ist, dass es eine injektive Abbildung
$f:M\to N$ gibt.
Es gilt $N=g^{-1}(M)=\bigcup_{x\in M}g^{-1}(\{x\})$, wobei jede Faser
$g^{-1}(\{x\})=\{y\in N\mid g(y)=x\}$ wegen der Surjektivität von $g$
nichtleer ist. Setze $X:=\{g^{-1}(\{x\})\mid x\in M\}$.
Nach dem Auswahlaxiom existiert eine Auswahlfunktion $\widetilde f:X\to N$,
das heißt $\widetilde f(g^{-1}(\{x\}))\in g^{-1}(\{x\})$ für alle $x\in M$.
Wir definieren nun
\[
    f:M\longrightarrow N,\qquad x\longmapsto\widetilde f\bigl(g^{-1}(\{x\})\bigr).
\]
Für jedes $x\in M$ gilt dann
$(g\circ f)(x)=g\bigl(\widetilde f(g^{-1}(\{x\}))\bigr)=x$, also ist
$g\circ f=\operatorname{id}_M$. Insbesondere ist $g\circ f$ injektiv,
woraus die Injektivität von $f$ folgt.
\end{proof}

\begin{bemerkung}[Injektivität bei Kompositionen]
Allgemein gilt: Sind $f:X\to Y$ und $g:Y\to Z$ Abbildungen, sodass
$g\circ f:X\to Z$ injektiv ist, dann ist auch $f$ injektiv (denn aus
$f(x_1)=f(x_2)$ folgt $g(f(x_1))=g(f(x_2))$ und damit $x_1=x_2$).
Die Abbildung $g$ muss dabei nicht injektiv sein, wie das Beispiel
$\{1\}\to\N\to\{1\}$ mit $1\mapsto1$ und $n\mapsto1$ zeigt.
\end{bemerkung}

\begin{bemerkung}[Existenz ohne explizite Konstruktion]
Als weiteres Beispiel für die Wirkung des Auswahlaxioms bzw.\ des Lemmas
von Zorn besitzt $\R$ als $\Q$-Vektorraum eine Basis, aber niemand hat
jemals eine solche Basis explizit angegeben.
\end{bemerkung}

\begin{satz}[Transitivität]
Seien $M_1,M_2,M_3$ Mengen. Dann gilt:
\begin{enumerate}
    \item $|M_1|\le|M_2|$ und $|M_2|\le|M_3|\implies |M_1|\le|M_3|$.
    \item $|M_1|=|M_2|$ und $|M_2|=|M_3|\implies |M_1|=|M_3|$.
\end{enumerate}
\end{satz}
\begin{proof}
Dies gilt, da die Verkettung injektiver (bzw.\ bijektiver) Abbildungen
wieder injektiv (bzw.\ bijektiv) ist.
\end{proof}

\begin{satz}[Satz von Cantor]
Für jede Menge $M$ gilt $|M|<|\Pot(M)|$.
(Für eine endliche Menge mit $|M|=n$ ist entsprechend $|\Pot(M)|=2^n>n$.)
\end{satz}
\begin{proof}
Wir wissen bereits, dass $|M|\le|\Pot(M)|$ gilt. Zu zeigen ist, dass es
keine bijektive (genauer: nicht einmal eine surjektive) Abbildung
$f:M\to\Pot(M)$ gibt.
Angenommen doch, es gäbe eine surjektive Abbildung $f:M\to\Pot(M)$.
Wir definieren
\[
    N:=\{x\in M\mid x\notin f(x)\}\in\Pot(M).
\]
Da $f$ surjektiv ist, existiert ein $x_0\in M$ mit $f(x_0)=N$.
Wir fragen, ob $x_0\in N$ gilt:
\begin{itemize}
    \item Falls $x_0\in N=\{x\in M\mid x\notin f(x)\}$, so folgt
    $x_0\notin f(x_0)=N$ -- Widerspruch.
    \item Falls $x_0\notin N=\{x\in M\mid x\notin f(x)\}$, so folgt
    $x_0\in f(x_0)=N$ -- Widerspruch.
\end{itemize}
Damit kann keine solche Abbildung $f$ existieren.
\end{proof}

\begin{bemerkung}
Durch wiederholte Potenzmengenbildung $\N\rightsquigarrow\Pot(\N)
\rightsquigarrow\Pot(\Pot(\N))\rightsquigarrow\dots$ erhält man
unendlich viele paarweise verschiedene unendliche Mächtigkeiten.
\end{bemerkung}

\begin{satz}[Vergleichssatz von Hartogs]
Seien $M$ und $N$ Mengen. Dann gilt $|M|\le|N|$ oder $|N|\le|M|$.
\end{satz}
\begin{proof}
Wir können annehmen, dass $M,N\ne\emptyset$ sind (ist eine der Mengen
leer, existiert die leere injektive Abbildung in die andere Menge).
Zu zeigen ist, dass es eine injektive Abbildung $f:M\to N$ oder eine
injektive Abbildung $g:N\to M$ gibt.
Betrachte die Menge
\[
    X:=\{f:M'\to N\mid M'\subset M\text{ Teilmenge},\ f\text{ injektiv}\}
\]
und definiere darauf eine Relation $\le$ durch
\[
    (f:M'\to N)\le(\bar f:M''\to N)
    :\iff M'\subset M''\text{ und }\bar f|_{M'}=f.
\]
Dies ist eine Halbordnung auf $X$. Wir prüfen die Voraussetzung des
Lemmas von Zorn: Sei $Y\subset X$ eine totalgeordnete Teilmenge.
(Für $Y=\emptyset$ ist die leere Abbildung $\emptyset\to N$ eine obere
Schranke in $X$.) Für $Y\ne\emptyset$ definiere eine obere Schranke
auf $U:=\bigcup_{(f:M'\to N)\in Y}M'$ durch
\[
    h:U\longrightarrow N,\qquad x\longmapsto g(x),
\]
wobei $(g:M'\to N)\in Y$ mit $x\in M'$ gewählt ist.
Da $Y$ totalgeordnet ist, ist dies wohldefiniert (liegen $x\in M_1'$ und
$x\in M_2'$ vor, so ist eine der Abbildungen Einschränkung der anderen)
und injektiv (für je zwei $x_1,x_2\in U$ gibt es ein gemeinsames
$(g:M'\to N)\in Y$ mit $x_1,x_2\in M'$).
Somit besitzt jede totalgeordnete Teilmenge $Y\subset X$ eine obere Schranke
in $X$. Nach dem Lemma von Zorn besitzt $X$ ein maximales Element
$f:M'\to N$.

\emph{(Abschluss des Beweises:)} Gilt $M'=M$, so ist $f:M\to N$ injektiv
und damit $|M|\le|N|$. Gilt $M'\ne M$, so muss $f(M')=N$ sein --
andernfalls ließen sich $x_0\in M\setminus M'$ und $y_0\in N\setminus f(M')$
wählen, wodurch man $f$ durch $x_0\mapsto y_0$ echt zu einer injektiven
Abbildung auf $M'\cup\{x_0\}$ fortsetzen könnte, im Widerspruch zur
Maximalität von $f$. Also ist $f:M'\to N$ bijektiv, und die Komposition
der Umkehrabbildung $f^{-1}:N\to M'$ mit der Inklusion $M'\hookrightarrow M$
liefert eine injektive Abbildung $N\to M$, also $|N|\le|M|$.
\end{proof}


\end{document}
